\documentclass[../main.tex]{subfiles}
\graphicspath{{\subfix{../image/}}{\subfix{../figure/}}}

\begin{document}

\chapter{案例研究二：双曲守恒律}\label{chapter:ch07}

\section{问题的提法}\label{section:ch07-问题提法}

本章考察 Cauchy 问题
\begin{align*}
\begin{cases}
\partial_t u + \partial_x[\phi(u)] = 0 & \text{在 }\mathbb{R}\times(0,T)=:S,\\
u=u_0 & \text{在 }\mathbb{R}\times\{0\},
\end{cases}
\qquad(P)
\end{align*}
其中 $T>0$，$u_0$ 为 $\mathbb{R}$ 上的正有限 Radon 测度。关于 $u_0$ 为 $\mathbb{R}$ 上有限带号 Radon 测度的情形，我们请读者参见文献~\cite{ref19}。

\subsection{假设与预备注记}\label{subsection:ch07-假设}

假设
\begin{align*}
u_0\in\mathscr{R}_f^+(\mathbb{R}), \qquad(A0)
\end{align*}
\begin{align*}
\phi\in\operatorname{Lip}([0,\infty)),\qquad \lim_{y\to\infty}\frac{\phi(y)}{y}=0, \qquad(A1)
\end{align*}
（不失一般性可设 $\phi(0)=0$）。需指出，$\lim_{y\to\infty}\phi(y)/y=C_\phi\neq0$ 的一般情形可通过将 $x$ 替换为 $x-C_\phi t$ 来处理（细节见文献~\cite{ref18}）。由 (A1)（及 $\phi(0)=0$），存在 $M>0$ 使
\begin{align}
|\phi'(y)|\leqslant M,\qquad |\phi(y)|\leqslant My\quad \text{对每个 }y>0. \label{align:ch07-7.1}
\end{align}
我们还将使用下列更强的假设（见\cref{section:ch07-唯一性}）：
\begin{align*}
\text{奇异部分 }u_{0s}\text{ 为 Dirac 质量的有限叠加，}
\quad u_{0s}=\sum_{j=1}^p c_j\delta_{x_j}\ (x_1<x_2<\cdots<x_p;\ c_j>0,\ 1\leqslant j\leqslant p). \qquad(A0')
\end{align*}
为特定目的，我们再作如下假设（见\cref{section:ch07-存在性}）：
\begin{align*}
\begin{cases}
\phi\in C^1([0,\infty)),\text{ 且对每个 }y_1>0,\text{ 存在 }a,b\geqslant0,\ a+b>0,\\
\text{使 }\phi'\text{ 在 }[y_1-a,y_1+b]\text{ 上严格单调},
\end{cases}
\qquad(A2)
\end{align*}
\begin{align*}
\begin{cases}
\phi\in C^\infty([0,\infty));\ \lim_{y\to\infty}\phi(y)/y=0;\\
\text{存在 }\alpha\geqslant-1\text{ 与 }\beta\in\mathbb{R}\text{ 使对一切 }y\in[0,\infty),\\
\phi''(y)[\alpha\phi(y)+\beta]\leqslant-[\phi'(y)]^2<0.
\end{cases}
\qquad(A2')
\end{align*}

\begin{definition}[流的类]\label{definition:ch07-流类}
称流 $\phi$ 属于\emph{正则流类}，若它满足上述结构假设 (A2')；其中常数 $\alpha\geqslant-1$、$\beta\in\mathbb{R}$ 区分\cref{remark:ch07-7.1.1}所述两种情形。存在唯一性与正则性结果（\cref{theorem:ch07-存在唯一}）即在该流类上建立。
\end{definition}

\begin{remark}\label{remark:ch07-7.1.1}
由 (A2')，映射 $y\mapsto\phi''(y)[\alpha\phi(y)+\beta]$ 在 $[0,\infty)$ 上严格负且连续，故只有两种可能：要么 (a) $\alpha\phi+\beta>0$ 且 $\phi''<0$，要么 (b) $\alpha\phi+\beta<0$ 且 $\phi''>0$，均在 $[0,\infty)$ 上。在情形 (a)，由 $\phi''<0$ 与 $\lim_{y\to\infty}\phi'(y)=0$ 知 $0<\phi'\leqslant\phi'(0)$ 在 $[0,\infty)$ 上。类似地，在情形 (b) 显然有 $\phi'(0)\leqslant\phi'<0$ 在 $[0,\infty)$ 上。特别地，两种情形下 (A2') 均蕴含 (A2) 与 (A1)。此外，若 $\phi(0)=0$，则 $\alpha\phi+\beta>0$ 在 $[0,\infty)$ 上当且仅当 $\beta>0$。
\end{remark}

流 $\phi$ 为超线性（形如 $\phi(y)=y^p$，$p>1$）时的问题 (P) 已在文献~\cite{ref67}中研究，其中证明了非负熵解的存在唯一性（另见文献~\cite{ref33}）。按该文定义，对正的时刻解取 $L^1(\mathbb{R})$ 中的值，尽管 $u_0$ 是有限 Radon 测度。关于 $\phi$ 在无穷远处至多线性增长时 (P) 的有趣结果，见开创性论文文献~\cite{ref39}。

当 $\phi(y)=Cy$（$C\in\mathbb{R}$）时，问题 (P) 即线性输运方程的 Cauchy 问题
\begin{align*}
\begin{cases}
u_t+Cu_x=0 & \text{在 }S,\\
u=u_0 & \text{在 }\mathbb{R}\times\{0\},
\end{cases}
\end{align*}
其解平凡地为 $u_0$ 沿直线 $x=Ct+x_0$（$x_0\in\mathbb{R}$）的平移。特别地，解的奇异部分 $u_s(\cdot,t)$ 对 $t>0$ 非零当且仅当它在 $t=0$ 时非零。

自然要问：当 $\phi$ 为次线性时会发生什么，若干数学模型正是如此提示（见文献~\cite{ref50,ref88,ref89}）。此情形下求 (P) 解的自然空间为 $L^\infty_{w^*}(0,T;\mathscr{R}_f^+(\mathbb{R}))$，如下列定义所设。

\subsection{解的定义}\label{subsection:ch07-解的定义}

本节始终假设 (A0)--(A1) 成立。

\begin{definition}\label{definition:ch07-young测度解}
问题 (P) 的 Young 测度解是一对 $(u,\nu)$，满足：
\begin{enumerate}[(i)]
\item $u\in L^\infty_{w^*}(0,T;\mathscr{R}_f^+(\mathbb{R}))$，$\nu\in\mathcal{Y}(S;\mathbb{R})$；
\item 对几乎所有 $(x,t)\in S$，$\operatorname{supp}\nu_{(x,t)}\subseteq[0,\infty)$，且
\begin{align}
u_r(x,t)=\int_{[0,\infty)}y\,\mathrm{d}\nu_{(x,t)}(y), \label{align:ch07-7.2}
\end{align}
其中 $\{\nu_{(x,t)}\}_{(x,t)\in S}$ 为 $\nu$ 的离散分解；
\item 对每个 $\zeta\in C^1([0,T];C^1_c(\mathbb{R}))$、$\zeta(\cdot,T)=0$ 在 $\mathbb{R}$ 上，有
\begin{align}
\iint_S[u_r\partial_t\zeta+\phi^*\partial_x\zeta]\,\mathrm{d}x\,\mathrm{d}t
+\int_0^T\langle u_s(\cdot,t),\partial_t\zeta(\cdot,t)\rangle\,\mathrm{d}t
=-\langle u_0,\zeta(\cdot,0)\rangle, \label{align:ch07-7.3}
\end{align}
其中
\begin{align}
\phi^*(x,t):=\int_{[0,\infty)}\phi(y)\,\mathrm{d}\nu_{(x,t)}(y)
\quad\text{对几乎所有 }(x,t)\in S. \label{align:ch07-7.4}
\end{align}
\end{enumerate}
\end{definition}

置 $H_+(y):=\chi_{(0,\infty)}(y)$，$H_-(y):=-\chi_{(-\infty,0)}(y)$，$\operatorname{sgn}(y):=H_+(y)+H_-(y)$（$y\in\mathbb{R}$）。

\begin{definition}\label{definition:ch07-young熵解}
问题 (P) 的 Young 测度熵解是一个 Young 测度解，满足对上述任意 $\zeta$、$\zeta\geqslant0$ 及 $l\in[0,\infty)$，
\begin{align}
\iint_S[E_l^*\partial_t\zeta+F_l^*\partial_x\zeta]\,\mathrm{d}x\,\mathrm{d}t
+\int_0^T\langle u_s(\cdot,t),\partial_t\zeta(\cdot,t)\rangle\,\mathrm{d}t
\geqslant-\int_{\mathbb{R}}E_l(u_{0r}(x))\zeta(x,0)\,\mathrm{d}x
-\langle u_{0s},\zeta(\cdot,0)\rangle, \label{align:ch07-7.5}
\end{align}
其中 $E_l(y)=|y-l|$，$F_l(y)=\operatorname{sgn}(y-l)[\phi(y)-\phi(l)]$（$y\in[0,\infty)$），且
\begin{align*}
E_l^*:=\int_{[0,\infty)}E_l(y)\,\mathrm{d}\nu(y),\qquad
F_l^*:=\int_{[0,\infty)}F_l(y)\,\mathrm{d}\nu(y).
\end{align*}
(P) 的 Young 测度熵下解定义为：对上述任意 $\zeta$ 与 $l$ 要求\eqref{align:ch07-7.5}成立，但取 $E_l(y)=[y-l]_+$、$F_l(y)=H_+(y-l)[\phi(y)-\phi(l)]$。Young 测度熵上解类似定义。
\end{definition}

\begin{remark}\label{remark:ch07-7.1.2}
(i) 由\eqref{align:ch07-7.1}、\eqref{align:ch07-7.2}与\eqref{align:ch07-7.4}得
\begin{align}
|\phi^*(x,t)|\leqslant M\int_{[0,\infty)}y\,\mathrm{d}\nu_{(x,t)}(y)=Mu_r(x,t)
\quad\text{对几乎所有 }(x,t)\in S. \label{align:ch07-7.6}
\end{align}
因 $u_r\in L^\infty(0,T;L^1(\mathbb{R}))$，由\eqref{align:ch07-7.6}知 $\phi^*\in L^\infty(0,T;L^1(\mathbb{R}))$。
\end{remark}

对任意开集 $\Omega\subseteq\mathbb{R}$ 与 $\tau\in(0,T)$，置 $Q_\tau:=\Omega\times(0,\tau)$。

\begin{definition}\label{definition:ch07-熵解}
测度 $u\in L^\infty_{w^*}(0,T;\mathscr{R}_f^+(\Omega))$ 称为 (P) 在 $Q_\tau$ 中的解，若对任意 $\zeta\in C^1([0,\tau];C^1_c(\Omega))$、$\zeta(\cdot,\tau)=0$ 在 $\Omega$ 上，有
\begin{align}
\iint_{Q_\tau}[u_r\partial_t\zeta+\phi(u_r)\partial_x\zeta]\,\mathrm{d}x\,\mathrm{d}t
+\int_0^\tau\langle u_s(\cdot,t),\partial_t\zeta(\cdot,t)\rangle\,\mathrm{d}t
=-\langle u_0,\zeta(\cdot,0)\rangle. \label{align:ch07-7.7}
\end{align}
(P) 在 $S$ 中的解简称为 (P) 的解。
\end{definition}

\begin{definition}\label{definition:ch07-Qtau熵解}
(P) 在 $Q_\tau$ 中的解称为 $Q_\tau$ 中的熵解，若它满足熵不等式：对所有 $\zeta\in C^1([0,\tau];C^1_c(\Omega))$、$\zeta\geqslant0$、$\zeta(\cdot,\tau)=0$ 在 $\Omega$ 上，以及所有 $l\in[0,\infty)$，
\begin{align}
\iint_{Q_\tau}\{E_l(u_r)\partial_t\zeta+F_l(u_r)\partial_x\zeta\}\,\mathrm{d}x\,\mathrm{d}t
+\int_0^\tau\langle u_s(\cdot,t),\partial_t\zeta(\cdot,t)\rangle\,\mathrm{d}t
\geqslant-\int_\Omega E_l(u_{0r}(x))\zeta(x,0)\,\mathrm{d}x
-\langle u_{0s},\zeta(\cdot,0)\rangle. \label{align:ch07-7.8}
\end{align}
\end{definition}

\begin{definition}\label{definition:ch07-上下解}
(P) 在 $Q_\tau$ 中的熵下解与熵上解分别由要求下列不等式成立来定义：
\begin{align}
\iint_{Q_\tau}\{[u_r-l]_+\partial_t\zeta+H_+(u_r-l)[\phi(u_r)-\phi(l)]\partial_x\zeta\}\,\mathrm{d}x\,\mathrm{d}t
+\int_0^\tau\langle u_s(\cdot,t),\partial_t\zeta(\cdot,t)\rangle_\Omega\,\mathrm{d}t
\geqslant-\int_\Omega[u_{0r}-l]_+\zeta(x,0)\,\mathrm{d}x
-\langle u_{0s},\zeta(\cdot,0)\rangle_\Omega, \label{align:ch07-7.9}
\end{align}
以及
\begin{align}
\iint_{Q_\tau}\{[u_r-l]_-\partial_t\zeta+H_-(u_r-l)[\phi(u_r)-\phi(l)]\partial_x\zeta\}\,\mathrm{d}x\,\mathrm{d}t
\geqslant-\int_\Omega[u_{0r}-l]_-\zeta(x,0)\,\mathrm{d}x \label{align:ch07-7.10}
\end{align}
对上述所有 $\zeta$ 与 $l$ 成立。

易见每个熵解同时是 (P) 的熵下解与熵上解。
\end{definition}

\begin{remark}\label{remark:ch07-7.1.3}
(i) 等式\eqref{align:ch07-7.7}亦可写为
\begin{align*}
\langle u,\partial_t\zeta\rangle+\iint_S\phi(u_r)\partial_x\zeta\,\mathrm{d}x\,\mathrm{d}t=-\langle u_0,\zeta(\cdot,0)\rangle,
\end{align*}
表明 $u_t=-[\phi(u_r)]_x$ 在 $\mathcal{D}^*(S)$ 中。

(ii) 问题 (P) 的解也是 Young 测度解。此外由\eqref{align:ch07-7.1}知 $\phi(u_r)\in L^\infty(0,T;L^1(\mathbb{R}))$。熵解、熵下解与熵上解有类似注记。
\end{remark}

下列结果与\cref{proposition:ch6-1.1}类似，证明亦类似，故从略（见文献~\cite{ref17} 命题 3.5）。

\begin{proposition}\label{proposition:ch07-7.1.1}
设 (A1) 成立，$(u,\nu)$ 为 (P) 的 Young 测度解。则：
\begin{enumerate}[(i)]
\item 存在零测集 $F^*\subseteq(0,T)$，使对一切 $t\in(0,T)\setminus F^*$ 与 $\rho\in C^1_c(\mathbb{R})$，
\begin{align}
\langle u(\cdot,t),\rho\rangle-\langle u_0,\rho\rangle
=\int_0^t\int_{\mathbb{R}}\phi^*(x,s)\rho'(x)\,\mathrm{d}x\,\mathrm{d}s; \label{align:ch07-7.11}
\end{align}
\item 对一切 $\rho\in C_c(\mathbb{R})$，
\begin{align}
\operatorname*{ess\,lim}_{t\to0^+}\langle u(\cdot,t),\rho\rangle=\langle u_0,\rho\rangle, \label{align:ch07-7.12}
\end{align}
\begin{align}
\operatorname*{ess\,lim}_{t\to t_0}\langle u(\cdot,t),\rho\rangle=\langle u(\cdot,t_0),\rho\rangle
\quad\text{对几乎所有 }t_0\in(0,T). \label{align:ch07-7.13}
\end{align}
\end{enumerate}
\end{proposition}

\begin{remark}\label{remark:ch07-7.1.4}
如\cref{remark:ch06-6.1.2}，易见映射 $t\mapsto u(\cdot,t)$ 有一个对任意 $t\in[0,T]$ 都有定义的代表元，使
\begin{align}
\lim_{t\to t_0}\langle u(\cdot,t),\rho\rangle=\langle u(\cdot,t_0),\rho\rangle
\quad\text{对所有 }t_0\in[0,T]\text{ 与 }\rho\in C_c(\mathbb{R}). \label{align:ch07-7.14}
\end{align}
今后凡对 $t$ 在 $[0,T]$ 的某子区间上陈述映射 $t\mapsto u(\cdot,t)$ 的性质时，均指此连续代表元。
\end{remark}

如第 6 章所处理的抛物问题（见\cref{section:ch06-唯一性}），正则性信息对证明唯一性很重要。这方面的首个结果部分地使人想起不等式\eqref{align:ch6-6.33}，陈述如下。

\begin{proposition}\label{proposition:ch07-7.1.2}
设 (A0)--(A1) 成立。
\begin{enumerate}[(i)]
\item 设 $u$ 为 (P) 的 Young 测度熵解。则
\begin{align}
u_s(\cdot,t_2)\leqslant u_s(\cdot,t_1)\quad\text{在 }\mathscr{R}_f^+(\mathbb{R})\text{ 中对几乎所有 }0\leqslant t_1\leqslant t_2\leqslant T. \label{align:ch07-7.15}
\end{align}
特别地，
\begin{align}
u_s(\cdot,t)\leqslant u_{0s}\quad\text{在 }\mathscr{R}_f^+(\mathbb{R})\text{ 中对几乎所有 }t\in(0,T). \label{align:ch07-7.16}
\end{align}
\item 设 $u$ 为问题 (P) 的解。则质量守恒成立：
\begin{align}
\|u(\cdot,t)\|_{\mathscr{R}_f(\mathbb{R})}=\|u_0\|_{\mathscr{R}_f(\mathbb{R})}
\quad\text{对几乎所有 }t\in(0,T). \label{align:ch07-7.17}
\end{align}
\end{enumerate}
\end{proposition}

\section{唯一性}\label{section:ch07-唯一性}

本节考虑 $\phi$ 有界的情形（$\phi$ 无界时的唯一性结果见\cref{remark:ch07-7.3.1}）。

对 (P) 的解 $u$，定义等待时间 $t_0\in[0,T]$：
\begin{align}
t_0:=\inf\{\tau\in(0,T]\mid u_s(\cdot,t)=0,\ u_r(\cdot,t)\in L^\infty(\mathbb{R})
\text{ 对几乎所有 }t\in(\tau,T)\}. \label{align:ch07-7.18}
\end{align}

\begin{proposition}\label{proposition:ch07-7.2.1}
设 $\phi$ 有界，(A0)--(A1) 成立。设对某 $x_0\in\mathbb{R}$ 有 $u_{0s}(\{x_0\})>0$，$u$ 为问题 (P) 的解。则由\eqref{align:ch07-7.18}定义的等待时间 $t_0$ 满足
\begin{align}
t_0\geqslant\min\left\{T,\frac{u_{0s}(\{x_0\})}{2\|\phi\|_\infty}\right\}>0. \label{align:ch07-7.19}
\end{align}
\end{proposition}

\begin{remark}\label{remark:ch07-7.2.1}
联系等式\eqref{align:ch07-7.12}，注意若 $u_{0s}\neq0$ 且等待时间 $t_0=0$，则映射 $t\mapsto u(\cdot,t)$ 在 $t=0$ 处不可能在 $\mathscr{R}_f(\mathbb{R})$ 的强拓扑下连续。反之，等待时间为正时强连续性可能成立（见文献~\cite{ref18} 定理 4.1）。
\end{remark}

等待时间以 $\phi$ 与 $u_0$ 表出的上界由下列命题给出，其证明依赖密度 $u_r$ 的 Aronson--Bénilan 型估计，我们请读者参见文献~\cite{ref18} 定理 3.11(ii)。等待时间的精细估计另见文献~\cite{ref20}。

\begin{proposition}\label{proposition:ch07-7.2.2}
设 (A0') 与 (A2') 成立，$\phi$ 有界，$\alpha>-1$，$|\beta|<\lim_{y\to\infty}|\phi(y)|=:\gamma$。设 $u$ 为由\cref{theorem:ch07-存在唯一}给出的问题 (P) 的熵解。则由\eqref{align:ch07-7.18}定义的等待时间 $t_0$ 满足
\begin{align}
t_0\leqslant\min\left\{T,\frac{(\alpha+1)\|u_0\|_{\mathscr{R}_f(\mathbb{R})}}{\gamma-|\beta|}\right\}. \label{align:ch07-7.20}
\end{align}
\end{proposition}

\begin{remark}\label{remark:ch07-7.2.2}
若 $\phi(y)=1-(1+y)^p$（$p<0$），则由\cref{proposition:intro-0.2.4}的证明可知\eqref{align:ch07-7.18}定义的等待时间为 $t_0=1$。故此情形下估计\eqref{align:ch07-7.20}是精确的，因 (A2') 在 $\alpha=p/(1-p)$、$\beta=-p/(1-p)$ 下成立，从而
\begin{align*}
\frac{(\alpha+1)\|\delta_0\|_{\mathscr{R}_f(\mathbb{R})}}{\gamma-|\beta|}
=\frac{(p/(1-p)+1)\|\delta_0\|_{\mathscr{R}_f(\mathbb{R})}}{1+p/(1-p)}=1.
\end{align*}
\end{remark}

借助\cref{proposition:ch07-7.1.2}与\cref{proposition:ch07-7.2.1}（其证明在\cref{section:ch07-正则证明}中给出），现在可处理唯一性。设 (A0') 与 (A1) 成立（后者中 $\phi$ 有界）。则 $u_{0s}=\sum_{j=1}^p c_j\delta_{x_j}$，置
\begin{align*}
I_1:=(-\infty,x_1),\qquad I_j:=(x_{j-1},x_j)\ (j=2,\dots,p),\qquad I_{p+1}:=(x_p,\infty),
\end{align*}
及 $S_j:=I_j\times(0,T)$（$j=1,\dots,p+1$）。由\cref{proposition:ch07-7.2.1}，若 $u$ 为问题 (P) 的解，则
\begin{align*}
\forall x_j\ \exists t_j\in(0,T]\ \text{使得}\ 
\begin{cases}
u_s(\cdot,t)(\{x_j\})>0 & \text{对几乎所有 }t\in[0,t_j),\\
u_s(\cdot,t)(\{x_j\})=0 & \text{对几乎所有 }t\in(t_j,T).
\end{cases}
\end{align*}
于是可陈述如下定义。

\begin{definition}\label{definition:ch07-相容条件}
设 $\phi$ 有界，(A0')--(A1) 成立，$j=1,\dots,p$，$\tau\in(0,t_j]$。(P) 的熵解称为在 $[0,\tau]$ 上于 $x_j$ 处满足相容条件，若对所有 $\beta\in C^1_c(0,\tau)$、$\beta\geqslant0$ 及 $l\in[0,\infty)$，
\begin{align*}
\pm\operatorname*{ess\,lim}_{x\to x_j^\pm}\int_0^\tau H_-(u_r(x,t)-l)[\phi(u_r(x,t))-\phi(l)]\beta(t)\,\mathrm{d}t\leqslant0. \qquad(C^\pm)
\end{align*}
\end{definition}

下列引理表明\cref{definition:ch07-相容条件}是适定的。

\begin{lemma}\label{lemma:ch07-7.2.3}
设 $u$ 为 (P) 的熵上解，$\beta\in C^1_c(0,T)$，$\beta\geqslant0$。则：
\begin{enumerate}[(i)]
\item 对每个 $l\in[0,\infty)$，函数
\begin{align}
x\mapsto-\int_0^T H_-(u_r(x,t)-l)[\phi(u_r(x,t))-\phi(l)]\beta(t)\,\mathrm{d}t+lT\|\beta'\|_\infty x \label{align:ch07-7.21}
\end{align}
的分布导数非负；
\item 对一切 $x_0\in\mathbb{R}$ 与 $l\in[0,\infty)$，存在有限极限
\begin{align}
\operatorname*{ess\,lim}_{x\to x_0^\pm}\int_0^T H_-(u_r(x,t)-l)[\phi(u_r(x,t))-\phi(l)]\beta(t)\,\mathrm{d}t. \label{align:ch07-7.22}
\end{align}
\end{enumerate}
\end{lemma}

\begin{proof}
设 $\alpha\in C^1_c(\mathbb{R})$，$\alpha\geqslant0$。在\eqref{align:ch07-7.10}中取 $Q_\tau=S$、$\zeta(x,t)=\alpha(x)\beta(t)$，得
\begin{align*}
\iint_S\{[u_r(x,t)-l]_-\alpha(x)\beta'(t)+H_-(u_r(x,t)-l)[\phi(u_r(x,t))-\phi(l)]\alpha'(x)\beta(t)\}\,\mathrm{d}x\,\mathrm{d}t\geqslant0.
\end{align*}
因 $0\leqslant[u_r-l]_-\leqslant l$，由上式得
\begin{align*}
-\int_{\mathbb{R}}\left(\int_0^T H_-(u_r(x,t)-l)[\phi(u_r(x,t))-\phi(l)]\beta(t)\,\mathrm{d}t\right)\alpha'(x)\,\mathrm{d}x
\leqslant lT\|\beta'\|_\infty\int_{\mathbb{R}}\alpha(x)\,\mathrm{d}x,
\end{align*}
由此得 (i)。

故函数\eqref{align:ch07-7.21}的分布导数为 Radon 测度。显然，$L^1_{\mathrm{loc}}(\mathbb{R})$ 中函数
\begin{align*}
\widetilde H(x):=-\int_0^T H_-(u_r(x,t)-l)[\phi(u_r(x,t))-\phi(l)]\beta(t)\,\mathrm{d}t
\end{align*}
的分布导数（记为 $\mu$）亦然。固定 $\bar x\in\mathbb{R}$，当 $x\geqslant\bar x$ 时置 $f_\mu(x):=\mu((\bar x,x])$，当 $x<\bar x$ 时置 $f_\mu(x):=-\mu((x,\bar x])$。则 $f_\mu$ 右连续，且在每个紧集 $K\subset\mathbb{R}$ 上（至多差一个可能依赖于 $K$ 的常数）与 $\widetilde H$ 几乎处处相等（例如见文献~\cite{ref05} 定理 3.28）。结论得证。
\end{proof}

现可陈述唯一性结果。

\begin{theorem}\label{theorem:ch07-唯一性}
设 $\phi$ 有界，(A0')--(A1) 成立。则问题 (P) 至多存在一个熵解，它属于 $C([0,T];\mathscr{R}_f^+(\mathbb{R}))$，且对所有 $j=1,\dots,p$ 在 $[0,t_j]$ 上于每个 $x_j$ 处满足相容条件。
\end{theorem}

为证\cref{theorem:ch07-唯一性}，需要若干引理。

\begin{lemma}\label{lemma:ch07-7.2.5}
设 $\phi$ 有界，(A0')--(A1) 成立，$u$ 为 (P) 的熵解。则对所有 $\beta\in C^1_c(0,T)$、$\beta\geqslant0$、$l\in[0,\infty)$ 及 $j=1,\dots,p$，存在有限极限
\begin{align}
\operatorname*{ess\,lim}_{x\to x_j^\pm}\int_0^T\operatorname{sgn}(u_r(x,t)-l)[\phi(u_r(x,t))-\phi(l)]\beta(t)\,\mathrm{d}t. \label{align:ch07-7.23}
\end{align}
\end{lemma}

\begin{proof}
只证右极限，左极限类似。固定 $j=1,\dots,p$。由\cref{proposition:ch07-7.1.2}，(P) 每个熵解的奇异部分关于时间非增，故由 (A0') 知对所有 $t\in[0,T]$ 有 $u_s(\cdot,t)(I_{j+1})=0$。设 $\alpha\in C^1_c(I_{j+1})$，$\alpha\geqslant0$。在\eqref{align:ch07-7.8}中取 $Q_\tau=S$、$\zeta(x,t)=\alpha(x)\beta(t)$，得
\begin{align*}
\iint_{S_{j+1}}\{|u_r(x,t)-l|\alpha(x)\beta'(t)+\operatorname{sgn}(u_r(x,t)-l)[\phi(u_r(x,t))-\phi(l)]\alpha'(x)\beta(t)\}\,\mathrm{d}x\,\mathrm{d}t\geqslant0.
\end{align*}
因 $0\leqslant|u_r-l|\leqslant u_r+l$，有
\begin{align*}
\begin{aligned}
-\int_{I_{j+1}}\alpha'(x)\bigg(&\int_0^T\operatorname{sgn}(u_r(x,t)-l)[\phi(u_r(x,t))-\phi(l)]\beta(t)\,\mathrm{d}t\bigg)\mathrm{d}x\\
&\leqslant\|\beta'\|_\infty\int_{I_{j+1}}\alpha(x)\left(\int_0^T u_r(x,t)\,\mathrm{d}t+lT\right)\mathrm{d}x.
\end{aligned}
\end{align*}
此不等式表明映射
\begin{align*}
x\mapsto-\int_0^T\operatorname{sgn}(u_r(x,t)-l)[\phi(u_r(x,t))-\phi(l)]\beta(t)\,\mathrm{d}t
+\|\beta'\|_\infty\left(\int_0^T\int_{x_j}^x u_r(y,t)\,\mathrm{d}y\,\mathrm{d}t+lTx\right)
\end{align*}
在 $I_{j+1}$ 中非负。如\cref{lemma:ch07-7.2.3}证明中那样论证即得结论。
\end{proof}

\begin{lemma}\label{lemma:ch07-7.2.6}
设 (A0')--(A1') 成立，$u$ 为 (P) 的熵解。则对每个 $j=1,\dots,p$：
\begin{enumerate}[(i)]
\item 存在 $h_j^-,h_j^+\in L^\infty(0,T)$，使对所有 $\beta\in C^1_c(0,T)$，
\begin{align}
\operatorname*{ess\,lim}_{x\to x_j^\pm}\int_0^T\phi(u_r(x,t))\beta(t)\,\mathrm{d}t
=\int_0^T h_j^\pm(t)\beta(t)\,\mathrm{d}t; \label{align:ch07-7.24}
\end{align}
\item 若 $u$ 在 $[0,\tau]$ 上于 $x_j$ 处满足相容条件 (C$^\pm$)，则
\begin{align}
h_j^-\leqslant\liminf_{l\to\infty}\phi(l)\leqslant\limsup_{l\to\infty}\phi(l)\leqslant h_j^+
\quad\text{在 }(0,\tau)\text{ 中几乎处处}. \label{align:ch07-7.25}
\end{align}
\end{enumerate}
\end{lemma}

\begin{remark}\label{remark:ch07-7.2.3}
由标准稠密性论证，从\eqref{align:ch07-7.24}得
\begin{align}
\operatorname*{ess\,lim}_{x\to x_j^\pm}\int_0^T\phi(u_r(x,t))\zeta(x,t)\,\mathrm{d}t
=\int_0^T h_j^\pm(t)\zeta(x_j,t)\,\mathrm{d}t \label{align:ch07-7.26}
\end{align}
对每个 $\zeta\in L^1(0,T;C_c(U_j))$ 成立，其中 $x_j\in U_j\subseteq\mathbb{R}$，$U_j$ 开。
\end{remark}

\begin{proof}[证明\cref{lemma:ch07-7.2.6}(i)]
只证右极限。因对 $y\in\mathbb{R}$ 有 $\operatorname{sgn}y=1+2H_-(y)$，由\eqref{align:ch07-7.22}--\eqref{align:ch07-7.23}，\eqref{align:ch07-7.24}左端极限存在且有限。另一方面，对每个收敛于 $x_j^+$ 的序列 $\{x_n\}$，序列 $\{\phi(u_r(x_n,\cdot))\}$ 在 $L^\infty(0,T)$ 中有界，故存在子列 $\{x_{n_k}\}\subseteq\{x_n\}$ 与函数 $h_j^+\in L^\infty(0,T)$，使
\begin{align*}
\{\phi(u_r(x_{n_k},\cdot))\}\stackrel{*}{\rightharpoonup}h_j^+\quad\text{在 }L^\infty(0,T).
\end{align*}
因对所有 $\beta\in C^1_c(0,T)$，\eqref{align:ch07-7.24}左端极限存在，由可分性论证易见 $h_j^+$ 与序列 $\{x_n$ 的选取无关。故 (i) 成立。

(ii) 因 $u$ 是 $I_{j+1}\times(0,\tau)$ 中 (P) 的解且 $u_s\llcorner(I_{j+1}\times(0,\tau))=0$，由\eqref{align:ch07-7.7}，对所有 $l\in[0,\infty)$ 与 $\zeta\in C^1([0,\tau];C^1_c(I_{j+1}))$、$\zeta(\cdot,\tau)=0$ 在 $I_{j+1}$ 上，有
\begin{align*}
\int_0^\tau\int_{I_{j+1}}\{(u_r-l)\zeta_t+[\phi(u_r)-\phi(l)]\zeta_x\}\,\mathrm{d}x\,\mathrm{d}t
=-\int_{I_{j+1}}[u_{0r}(x)-l]\zeta(x,0)\,\mathrm{d}x.
\end{align*}
令
\begin{align}
\eta_\epsilon(x):=\frac{2(x-x_j)-\epsilon}{\epsilon}\chi_{[x_j+\epsilon/2,\,x_j+\epsilon]}(x)+\chi_{(x_j+\epsilon,\,x_{j+1}]}(x)
\quad(x\in I_{j+1}), \label{align:ch07-7.27}
\end{align}
并取 $\zeta\in C^1([0,\tau];C^1_c([x_j,x_{j+1})))$ 使 $\zeta(\cdot,\tau)=0$ 在 $I_{j+1}$ 上（$j=p$ 时 $x_{j+1}=\infty$）。由标准论证可在上式中取 $\zeta=\zeta\eta_\epsilon$，得
\begin{align*}
\begin{aligned}
\int_0^\tau\int_{I_{j+1}}&\{(u_r-l)\zeta_t\eta_\epsilon+[\phi(u_r)-\phi(l)]\zeta_x\eta_\epsilon\}\,\mathrm{d}x\,\mathrm{d}t
+\int_{I_{j+1}}[u_{0r}(x)-l]\zeta(x,0)\eta_\epsilon(x)\,\mathrm{d}x\\
&=-\frac{2}{\epsilon}\int_0^\tau\int_{x_j+\epsilon/2}^{x_j+\epsilon}[\phi(u_r)-\phi(l)]\zeta\,\mathrm{d}x\,\mathrm{d}t.
\end{aligned}
\end{align*}
令 $\epsilon\to0^+$，由上式即得（见\eqref{align:ch07-7.26}）
\begin{align}
\begin{aligned}
\int_0^\tau\int_{I_{j+1}}&\{(u_r-l)\zeta_t+[\phi(u_r)-\phi(l)]\zeta_x\}\,\mathrm{d}x\,\mathrm{d}t
+\int_{I_{j+1}}[u_{0r}(x)-l]\zeta(x,0)\,\mathrm{d}x\\
&=-\operatorname*{ess\,lim}_{x\to x_j^+}\int_0^\tau[\phi(u_r(x,t))-\phi(l)]\zeta(x,t)\,\mathrm{d}t
=-\int_0^\tau[h_j^+(t)-\phi(l)]\zeta(x_j,t)\,\mathrm{d}t.
\end{aligned} \label{align:ch07-7.28}
\end{align}
因 $u$ 是 $I_{j+1}\times(0,\tau)$ 中 (P) 的熵解，如前论证得
\begin{align*}
\begin{aligned}
\int_0^\tau\int_{I_{j+1}}&\{|u_r-l|\zeta_t+\operatorname{sgn}(u_r-l)[\phi(u_r)-\phi(l)]\zeta_x\}\,\mathrm{d}x\,\mathrm{d}t
+\int_{I_{j+1}}|u_{0r}(x)-l|\zeta(x,0)\,\mathrm{d}x\\
&\geqslant-\operatorname*{ess\,lim}_{x\to x_j^+}\int_0^\tau\operatorname{sgn}(u_r(x,t)-l)[\phi(u_r(x,t))-\phi(l)]\zeta(x,t)\,\mathrm{d}t
\end{aligned}
\end{align*}
对上述一切 $\zeta$、$\zeta\geqslant0$ 成立。取 $\zeta(x,t)=\alpha(x)\beta(t)$，其中 $\alpha\in C^1_c([x_j,x_{j+1}))$、$\alpha\geqslant0$，$\beta\in C^1([0,\tau])$、$\beta\geqslant0$、$\beta(\tau)=0$。由相容条件 (C$^\pm$) 得
\begin{align}
\begin{aligned}
\int_0^\tau\int_{I_{j+1}}&\{|u_r-l|\zeta_t+\operatorname{sgn}(u_r-l)[\phi(u_r)-\phi(l)]\zeta_x\}\,\mathrm{d}x\,\mathrm{d}t
+\int_{I_{j+1}}|u_{0r}(x)-l|\zeta(x,0)\,\mathrm{d}x\\
&+\operatorname*{ess\,lim}_{x\to x_j^+}\int_0^\tau[\phi(u_r(x,t))-\phi(l)]\zeta(x,t)\,\mathrm{d}t
\geqslant-2\operatorname*{ess\,lim}_{x\to x_j^+}\int_0^\tau H_-(u_r(x,t)-l)[\phi(u_r(x,t))-\phi(l)]\zeta(x,t)\,\mathrm{d}t\\
&=-2\alpha(x_j)\operatorname*{ess\,lim}_{x\to x_j^+}\int_0^\tau H_-(u_r(x,t)-l)[\phi(u_r(x,t))-\phi(l)]\beta(t)\,\mathrm{d}t\geqslant0,
\end{aligned} \label{align:ch07-7.29}
\end{align}
因 $\operatorname{sgn}(y)=1+2H_-(y)$。由\eqref{align:ch07-7.28}与\eqref{align:ch07-7.29}得
\begin{align*}
\begin{aligned}
\int_0^\tau\int_{I_{j+1}}&\{[u_r-l]_+\zeta_t+H_+(u_r-l)[\phi(u_r)-\phi(l)]\zeta_x\}\,\mathrm{d}x\,\mathrm{d}t
+\int_{I_{j+1}}[u_{0r}(x)-l]_+\zeta(x,0)\,\mathrm{d}x\\
&\geqslant-\int_0^\tau[h_j^+(t)-\phi(l)]\zeta(x_j,t)\,\mathrm{d}t.
\end{aligned}
\end{align*}
令 $l\to\infty$，得
\begin{align*}
\liminf_{l\to\infty}\int_0^\tau[h_j^+(t)-\phi(l)]\zeta(x_j,t)\,\mathrm{d}t
=\int_0^\tau\left[h_j^+(t)-\limsup_{l\to\infty}\phi(l)\right]\zeta(x_j,t)\,\mathrm{d}t\geqslant0,
\end{align*}
由 $\zeta$ 的任意性得\eqref{align:ch07-7.25}中后一不等式。

将 $I_{j+1}\times(0,\tau)$ 换为 $I_j\times(0,\tau)$，类似\eqref{align:ch07-7.28}与\eqref{align:ch07-7.29}可得
\begin{align}
\begin{aligned}
\int_0^\tau\int_{I_j}&\{(u_r-l)\zeta_t+[\phi(u_r)-\phi(l)]\zeta_x\}\,\mathrm{d}x\,\mathrm{d}t
+\int_{I_j}[u_{0r}(x)-l]\zeta(x,0)\,\mathrm{d}x\\
&=\operatorname*{ess\,lim}_{x\to x_j^-}\int_0^\tau[\phi(u_r(x,t))-\phi(l)]\zeta(x,t)\,\mathrm{d}t
=\int_0^\tau[h_j^-(t)-\phi(l)]\zeta(x_j,t)\,\mathrm{d}t,
\end{aligned} \label{align:ch07-7.30}
\end{align}
\begin{align}
\begin{aligned}
\int_0^\tau\int_{I_j}&\{|u_r-l|\zeta_t+\operatorname{sgn}(u_r-l)[\phi(u_r)-\phi(l)]\zeta_x\}\,\mathrm{d}x\,\mathrm{d}t
+\int_{I_j}|u_{0r}(x)-l|\zeta(x,0)\,\mathrm{d}x\\
&-\operatorname*{ess\,lim}_{x\to x_j^-}\int_0^\tau[\phi(u_r(x,t))-\phi(l)]\zeta(x,t)\,\mathrm{d}t\geqslant0,
\end{aligned} \label{align:ch07-7.31}
\end{align}
从而
\begin{align*}
\begin{aligned}
\int_0^\tau\int_{I_j}&\{[u_r-l]_+\zeta_t+H_+(u_r-l)[\phi(u_r)-\phi(l)]\zeta_x\}\,\mathrm{d}x\,\mathrm{d}t
+\int_{I_j}[u_{0r}(x)-l]_+\zeta(x,0)\,\mathrm{d}x\\
&\geqslant\int_0^\tau[h_j^-(t)-\phi(l)]\zeta(x_j,t)\,\mathrm{d}t
\end{aligned}
\end{align*}
且
\begin{align*}
\limsup_{l\to\infty}\int_0^\tau[h_j^-(t)-\phi(l)]\zeta(x_j,t)\,\mathrm{d}t
=\int_0^\tau\left[h_j^-(t)-\liminf_{l\to\infty}\phi(l)\right]\zeta(x_j,t)\,\mathrm{d}t\leqslant0.
\end{align*}
由 $\zeta$ 任意得\eqref{align:ch07-7.25}中前一不等式。
\end{proof}

\begin{remark}\label{remark:ch07-7.2.4}
由标准稠密性论证与\eqref{align:ch07-7.26}，从\eqref{align:ch07-7.31}得：若 $\zeta\in C^1([0,\tau];C^1_c((-\infty,x_1]))$、$\zeta\geqslant0$、$\zeta(\cdot,\tau)=0$ 在 $(-\infty,x_1]$ 上，则
\begin{align}
\begin{aligned}
\int_0^\tau\int_{I_1}&\{|u_r-l|\zeta_t+\operatorname{sgn}(u_r-l)[\phi(u_r)-\phi(l)]\zeta_x\}\,\mathrm{d}x\,\mathrm{d}t
+\int_{I_1}|u_{0r}(x)-l|\zeta(x,0)\,\mathrm{d}x\\
&\geqslant\int_0^\tau[h_1^-(t)-\phi(l)]\zeta(x_1,t)\,\mathrm{d}t,
\end{aligned} \label{align:ch07-7.32}
\end{align}
并从\eqref{align:ch07-7.29}得：对所有 $\zeta\in C^1([0,\tau];C^1_c([x_p,\infty)))$、$\zeta\geqslant0$、$\zeta(\cdot,\tau)=0$ 在 $[x_p,\infty)$ 上，
\begin{align*}
\begin{aligned}
\int_0^\tau\int_{I_{p+1}}&\{|u_r-l|\zeta_t+\operatorname{sgn}(u_r-l)[\phi(u_r)-\phi(l)]\zeta_x\}\,\mathrm{d}x\,\mathrm{d}t
+\int_{I_{p+1}}|u_{0r}(x)-l|\zeta(x,0)\,\mathrm{d}x\\
&\geqslant-\int_0^\tau[h_p^+(t)-\phi(l)]\zeta(x_p,t)\,\mathrm{d}t.
\end{aligned}
\end{align*}
此外，将\eqref{align:ch07-7.27}中 $\eta_\epsilon$ 换为
\begin{align*}
\frac{2(x-x_j)-\epsilon}{\epsilon}\chi_{[x_j+\epsilon/2,\,x_j+\epsilon]}+\chi_{[x_j+\epsilon,\,x_{j+1}-\epsilon]}+\frac{2(x_{j+1}-x)-\epsilon}{\epsilon}\chi_{[x_{j+1}-\epsilon,\,x_{j+1}-\epsilon/2]},
\end{align*}
如\cref{lemma:ch07-7.2.6}证明那样论证，可得对任意 $j=1,\dots,p-1$，
\begin{align*}
\begin{aligned}
\int_0^\tau\int_{I_{j+1}}&\{|u_r-l|\zeta_t+\operatorname{sgn}(u_r-l)[\phi(u_r)-\phi(l)]\zeta_x\}\,\mathrm{d}x\,\mathrm{d}t
+\int_{I_{j+1}}|u_{0r}(x)-l|\zeta(x,0)\,\mathrm{d}x\\
&\geqslant-\int_0^\tau[h_j^+(t)-\phi(l)]\zeta(x_j,t)\,\mathrm{d}t
+\int_0^\tau[h_{j+1}^-(t)-\phi(l)]\zeta(x_{j+1},t)\,\mathrm{d}t
\end{aligned}
\end{align*}
对所有 $\zeta\in C^1([0,\tau];C^1_c([x_j,x_{j+1}])$、$\zeta\geqslant0$、$\zeta(\cdot,\tau)=0$ 在 $[x_j,x_{j+1}]$ 上成立。
\end{remark}

\begin{remark}\label{remark:ch07-7.2.5}
给出下列不等式，它们对所有 $\zeta\in C^1([0,\tau];C^1_c((-\infty,x_1]))$、$\zeta\geqslant0$、$\zeta(\cdot,\tau)=0$ 在 $(-\infty,x_1]$ 上成立：
\begin{align}
\begin{aligned}
\int_0^\tau\int_{I_1}&\{[u_r-l]_+\zeta_t+H_+(u_r-l)[\phi(u_r)-\phi(l)]\zeta_x\}\,\mathrm{d}x\,\mathrm{d}t
+\int_{I_1}[u_{0r}(x)-l]_+\zeta(x,0)\,\mathrm{d}x\\
&\geqslant\int_0^\tau[h_1^-(t)-\phi(l)]\zeta(x_1,t)\,\mathrm{d}t,
\end{aligned} \label{align:ch07-7.33}
\end{align}
\begin{align}
\begin{aligned}
\int_0^\tau\int_{I_1}&\{[u_r-l]_-\zeta_t+H_-(u_r-l)[\phi(u_r)-\phi(l)]\zeta_x\}\,\mathrm{d}x\,\mathrm{d}t
+\int_{I_1}[u_{0r}(x)-l]_-\zeta(x,0)\,\mathrm{d}x\geqslant0.
\end{aligned} \label{align:ch07-7.34}
\end{align}
其证明由\eqref{align:ch07-7.9}与\eqref{align:ch07-7.10}（而非\eqref{align:ch07-7.8}）出发，类似\eqref{align:ch07-7.32}而得。$j=2,\dots,p+1$ 的 $S_j$ 中有类似不等式。
\end{remark}

现可证明\cref{theorem:ch07-唯一性}。

\begin{proof}[证明\cref{theorem:ch07-唯一性}]
设 $u,v\in C([0,T];\mathscr{R}_f^+(\Omega))$ 为满足在每个 $x_j$ 于 $[0,t_j]$ 处相容条件的熵解，令
\begin{align}
\tau:=\min\{t_u,t_v\},\quad
\text{其中}\ 
\begin{cases}
t_u:=\sup\{t\in[0,T)\mid\operatorname{supp}u_s(\cdot,t)=\operatorname{supp}u_{0s}\},\\
t_v:=\sup\{t\in[0,T)\mid\operatorname{supp}v_s(\cdot,t)=\operatorname{supp}u_{0s}\}.
\end{cases} \label{align:ch07-7.35}
\end{align}
只需证
\begin{align}
u=v\quad\text{在 }\mathscr{R}_f(S_\tau), \label{align:ch07-7.36}
\end{align}
因若 $\tau<T$，标准迭代程序即可得结论。显然，若证得
\begin{align}
u_r=v_r\quad\text{在 }S_\tau\text{ 中几乎处处}, \label{align:ch07-7.37}
\end{align}
则\eqref{align:ch07-7.36}成立。事实上，由\eqref{align:ch07-7.7}与\eqref{align:ch07-7.37}得
\begin{align*}
\int_0^\tau\langle u_s(\cdot,t)-v_s(\cdot,t),\partial_t\zeta(\cdot,t)\rangle\,\mathrm{d}t
=-\iint_{S_\tau}\{(u_r-v_r)\partial_t\zeta+[\phi(u_r)-\phi(u_r)]\partial_x\zeta\}\,\mathrm{d}x\,\mathrm{d}t=0
\end{align*}
对所有 $\zeta\in C^1([0,\tau];C^1_c(\mathbb{R}))$、$\zeta(\cdot,\tau)=0$ 成立。故对几乎所有 $t\in(0,\tau)$ 与所有 $\alpha\in C^1_c(\mathbb{R})$，$\langle u_s(\cdot,t)-v_s(\cdot,t),\alpha\rangle=0$。因此 $u_s=v_s$ 在 $L^\infty(0,\tau;\mathscr{R}_f(\mathbb{R}))$，结合\eqref{align:ch07-7.37}得\eqref{align:ch07-7.36}。

余下需证\eqref{align:ch07-7.37}，等价于证对任意 $j=1,\dots,p+1$，
\begin{align}
u_r=v_r\quad\text{在 }I_j\times(0,\tau)\text{ 中几乎处处}. \label{align:ch07-7.38}
\end{align}
只证 $j=1$，其余情形类似。置 $Q_1:=(-\infty,x_1]\times(0,\tau)$。我们应用适配边值问题的 Kružkov 双变量法（见文献~\cite{ref68,ref75,ref91}）。设 $\xi=\xi(x,t,y,s)\geqslant0$ 定义于 $Q_1\times Q_1$，使对每个 $(y,s)\in Q_1$ 有 $\xi(\cdot,\cdot,y,s)\in C^1_c(Q_1)$，对每个 $(x,t)\in Q_1$ 有 $\xi(x,t,\cdot,\cdot)\in C^1_c(Q_1)$。由\eqref{align:ch07-7.32}得
\begin{align*}
\begin{aligned}
\iint_{Q_1}\{&\operatorname{sgn}(u_r(x,t)-v(y,s))[\phi(u_r(x,t))-\phi(v_r(y,s))]\xi_x(x,t,y,s)\\
&+|u_r(x,t)-v_r(y,s)|\xi_t(x,t,y,s)\}\,\mathrm{d}x\,\mathrm{d}t
\geqslant\int_0^\tau[h_1^-(t)-\phi(v_r(y,s))]\xi(x_1,t,y,s)\,\mathrm{d}t,
\end{aligned}
\end{align*}
\begin{align*}
\begin{aligned}
\iint_{Q_1}\{&\operatorname{sgn}(u_r(x,t)-v(y,s))[\phi(u_r(x,t))-\phi(v_r(y,s))]\xi_y(x,t,y,s)\\
&+|u_r(x,t)-v_r(y,s)|\xi_s(x,t,y,s)\}\,\mathrm{d}y\,\mathrm{d}s
\geqslant\int_0^\tau[g_1^-(s)-\phi(u_r(x,t))]\xi(x,t,x_1,s)\,\mathrm{d}s,
\end{aligned}
\end{align*}
其中由\cref{lemma:ch07-7.2.6}(i)，$g_j^\pm\in L^\infty(0,T)$ 对每个 $j=1,\dots,p$ 满足：若 $\beta\in C^1_c(0,T)$，
\begin{align}
\operatorname*{ess\,lim}_{x\to x_j^\pm}\int_0^T\phi(v_r(x,t))\beta(t)\,\mathrm{d}t
=\int_0^T g_j^\pm(t)\beta(t)\,\mathrm{d}t. \label{align:ch07-7.39}
\end{align}
设 $\rho_\epsilon$（$\epsilon>0$）为 $\mathbb{R}$ 上对称磨光核，在上式中取
\begin{align}
\xi(x,t,y,s)=\eta\left(\frac{x+y}{2},\frac{t+s}{2}\right)\rho_\epsilon(x-y)\rho_\epsilon(t-s) \label{align:ch07-7.40}
\end{align}
其中 $\eta\in C^1_c((-\infty,x_1]\times(0,\tau))$，$\eta\geqslant0$。于是得
\begin{align}
\begin{aligned}
\iint_{Q_1\times Q_1}&\rho_\epsilon(x-y)\rho_\epsilon(t-s)\Big\{|u_r(x,t)-v_r(y,s)|\eta_t\left(\frac{x+y}{2},\frac{t+s}{2}\right)\\
&+\operatorname{sgn}(u_r(x,t)-v_r(y,s))[\phi(u_r(x,t))-\phi(v_r(y,s))]\eta_x\left(\frac{x+y}{2},\frac{t+s}{2}\right)\Big\}\,\mathrm{d}x\,\mathrm{d}t\,\mathrm{d}y\,\mathrm{d}s\\
\geqslant\int_0^\tau\iint_{Q_1}&[g_1^-(s)-\phi(u_r(x,t))]\eta\left(\frac{x+x_1}{2},\frac{t+s}{2}\right)\rho_\epsilon(x_1-x)\rho_\epsilon(t-s)\,\mathrm{d}x\,\mathrm{d}t\,\mathrm{d}s\\
+&\int_0^\tau\iint_{Q_1}[h_1^-(t)-\phi(v_r(y,s))]\eta\left(\frac{x_1+y}{2},\frac{t+s}{2}\right)\rho_\epsilon(y-x_1)\rho_\epsilon(t-s)\,\mathrm{d}y\,\mathrm{d}s\,\mathrm{d}t.
\end{aligned} \label{align:ch07-7.41}
\end{align}
关于\eqref{align:ch07-7.41}右端，由磨光核熟知性质，当 $\epsilon\to0^+$ 时
\begin{align*}
\int_0^\tau\iint_{Q_1}g_1^-(s)\eta\left(\frac{x+x_1}{2},\frac{t+s}{2}\right)\rho_\epsilon(x_1-x)\rho_\epsilon(t-s)\,\mathrm{d}x\,\mathrm{d}t\,\mathrm{d}s
\to\frac12\int_0^\tau g_1^-(s)\eta(x_1,s)\,\mathrm{d}s,
\end{align*}
\begin{align*}
\int_0^\tau\iint_{Q_1}h_1^-(t)\eta\left(\frac{x_1+y}{2},\frac{t+s}{2}\right)\rho_\epsilon(y-x_1)\rho_\epsilon(t-s)\,\mathrm{d}y\,\mathrm{d}s\,\mathrm{d}t
\to\frac12\int_0^\tau h_1^-(t)\eta(x_1,t)\,\mathrm{d}t.
\end{align*}
此外，因 $\iint_{Q_1}\rho_\epsilon(x_1-x)\rho_\epsilon(t-s)\,\mathrm{d}x\,\mathrm{d}s=\frac12$（当 $\epsilon<\min\{t,\tau-t\}$），由 $\eta$ 的光滑性与等式\eqref{align:ch07-7.26}，差
\begin{align*}
\int_0^\tau\iint_{Q_1}\phi(u_r(x,t))\eta\left(\frac{x+x_1}{2},\frac{t+s}{2}\right)\rho_\epsilon(x_1-x)\rho_\epsilon(t-s)\,\mathrm{d}x\,\mathrm{d}t\,\mathrm{d}s
-\frac12\int_0^\tau h_1^-(t)\eta(x_1,t)\,\mathrm{d}t
\end{align*}
当 $\epsilon\to0^+$ 时趋于零。类似地，
\begin{align*}
\int_0^\tau\iint_{Q_1}\phi(v_r(y,s))\eta\left(\frac{x_1+y}{2},\frac{t+s}{2}\right)\rho_\epsilon(y-x_1)\rho_\epsilon(t-s)\,\mathrm{d}y\,\mathrm{d}s\,\mathrm{d}t
\to\frac12\int_0^\tau g_1^-(s)\eta(x_1,s)\,\mathrm{d}s.
\end{align*}
在\eqref{align:ch07-7.41}中令 $\epsilon\to0^+$，得对每个 $\eta\in C^1_c((-\infty,x_1]\times(0,\tau))$、$\eta\geqslant0$，
\begin{align}
\iint_{Q_1}\{|u_r(x,t)-v_r(x,t)|\eta_t(x,t)+\operatorname{sgn}(u_r-v_r)[\phi(u_r)-\phi(v_r)]\eta_x(x,t)\}\,\mathrm{d}x\,\mathrm{d}t\geqslant0. \label{align:ch07-7.42}
\end{align}
现固定 $t'$、$t''$ 使 $0<t'<t''<\tau$，并取 $x_0<x_1$。设 $\alpha_\delta=\alpha_\delta(x)$、$\beta_\vartheta=\beta_\vartheta(t)$ 为磨光核族，$0<\delta<1$，$0<\vartheta<\min\{t',\tau-t''\}$。在\eqref{align:ch07-7.42}中取
\begin{align*}
\eta(x,t)=\eta_{\delta,\vartheta}(x,t):=\int_{t-t''}^{t-t'}\beta_\vartheta(z)\,\mathrm{d}z\int_{\|\phi'\|_\infty(t-t'')+x_0}^{x_1+\delta}\alpha_\delta(x-y)\,\mathrm{d}y\quad(\delta>0)
\end{align*}
其中 $x\in(\|\phi'\|_\infty(t-t'')+x_0,x_1]$，$t\in(t'-\vartheta,t''+\vartheta)$（显然 $\eta_{\delta,\vartheta}\geqslant0$ 且属于 $C^\infty_c((-\infty,x_1]\times(0,\tau))$）。因 $x\in(-\infty,x_1]$ 时 $\alpha_\delta(x-x_1-\delta)=0$，
\begin{align*}
\begin{aligned}
\iint_{Q_1}&|u_r-v_r|[\beta_\vartheta(t-t')-\beta_\vartheta(t-t'')]\left(\int_{\|\phi'\|_\infty(t-t'')+x_0}^{x_1+\delta}\alpha_\delta(x-y)\,\mathrm{d}y\right)\mathrm{d}x\,\mathrm{d}t\\
-&\iint_{Q_1}\{\|\phi'\|_\infty|u_r-v_r|+\operatorname{sgn}(u_r-v_r)[\phi(u_r)-\phi(v_r)]\}\\
&\qquad\times\alpha_\delta(x-\|\phi'\|_\infty(t-t'')-x_0)\left(\int_{t-t''}^{t-t'}\beta_\vartheta(z)\,\mathrm{d}z\right)\mathrm{d}x\,\mathrm{d}t\geqslant0.
\end{aligned}
\end{align*}
因对一切 $y_1,y_2\geqslant0$，$\|\phi'\|_\infty|y_1-y_2|+\operatorname{sgn}(y_1-y_2)[\phi(y_1)-\phi(y_2)]\geqslant0$，故
\begin{align*}
\iint_{Q_1}|u_r-v_r|[\beta_\vartheta(t-t')-\beta_\vartheta(t-t'')]\left(\int_{\|\phi'\|_\infty(t-t'')+x_0}^{x_1+\delta}\alpha_\delta(x-y)\,\mathrm{d}y\right)\mathrm{d}x\,\mathrm{d}t\geqslant0.
\end{align*}
令 $\delta\to0^+$，由控制收敛定理得
\begin{align*}
\int_0^\tau\int_{\|\phi'\|_\infty(t-t'')+x_0}^{x_1}|u_r(x,t)-v_r(x,t)|[\beta_\vartheta(t-t')-\beta_\vartheta(t-t'')]\,\mathrm{d}x\,\mathrm{d}t\geqslant0,
\end{align*}
从而令 $\vartheta\to0^+$（回忆 $u,v\in C([0,T];\mathscr{R}_f^+(\Omega))$ 时 $u_r,v_r\in C([0,T];L^1(\mathbb{R}))$），得
\begin{align}
\int_{x_0}^{x_1}|u_r(x,t'')-v_r(x,t'')|\,\mathrm{d}x
\leqslant\int_{x_0-\|\phi'\|_\infty(t''-t')}^{x_1}|u_r(x,t')-v_r(x,t')|\,\mathrm{d}x. \label{align:ch07-7.43}
\end{align}
因 $u,v\in C([0,T];\mathscr{R}_f^+(\Omega))$，在\eqref{align:ch07-7.43}中令 $t'\to0^+$，对一切 $(x_0,t'')\in Q_1$ 得
\begin{align*}
\int_{x_0}^{x_1}|u_r(x,t'')-v_r(x,t'')|\,\mathrm{d}x
\leqslant\int_{x_0-\|\phi'\|_\infty t''}^{x_1}|u_r(x,0)-v_r(x,0)|\,\mathrm{d}x=0,
\end{align*}
因 $u_r(\cdot,0)=v_r(\cdot,0)=u_{0r}$。由 $t''\in(0,\tau)$ 任意得 $u_r=v_r$ 在 $Q_1$ 中。
\end{proof}

\begin{remark}\label{remark:ch07-7.2.6}
如\cref{remark:ch06-6.3.2}那样论证，易构造问题 (P) 不满足 (C$^\pm$)、从而不是所构造解（见\cref{definition:ch07-构造解}与\cref{theorem:ch07-存在唯一}）的熵解。故无相容条件时唯一性失效。
\end{remark}

\section{存在性与正则性结果}\label{section:ch07-存在性}

解的存在性通过一个逼近程序证明。考虑逼近问题
\begin{align*}
\begin{cases}
\partial_t u_n+\partial_x[\phi(u_n)]=0 & \text{在 }S,\\
u_n=u_{0n} & \text{在 }\mathbb{R}\times\{0\}
\end{cases}
\qquad(Pn)
\end{align*}
研究序列 $\{u_n\}$ 的极限点，我们将证明下列结果。

\begin{theorem}\label{theorem:ch07-7.3.1}
设 (A0)--(A1) 成立。则问题 (P) 有解 $u$，它是问题 (Pn) 的熵解序列 $\{u_n\}$ 的一个极限点。此外，$u$ 是 (P) 的 Young 测度熵解。
\end{theorem}

\begin{theorem}\label{theorem:ch07-7.3.2}
设 (A0)--(A2) 成立，$u$ 为由\cref{theorem:ch07-7.3.1}给出的 (P) 的解。则 $u$ 是 (P) 的熵解。
\end{theorem}

下列定义是\cref{definition:ch06-构造解}的对应物。

\begin{definition}\label{definition:ch07-构造解}
由\cref{theorem:ch07-7.3.1}与\cref{theorem:ch07-7.3.2}给出的 (P) 的解分别称为所构造解与所构造熵解。
\end{definition}

如\cref{section:ch06-存在正则}，(P) 的所构造熵解有若干重要性质。特别地，在附加条件 (A2')（它反过来蕴含 (A1) 与 (A2)）下，它们满足相容条件。这就是下列定理的内容。

\begin{theorem}\label{theorem:ch07-7.3.3}
设 (A0) 与 (A2') 成立，$u$ 为问题 (P) 的所构造熵解。则对每个 $\tau\in(0,T)$：
\begin{enumerate}[(i)]
\item 若 $\phi$ 有界，则 $u$ 在 $[0,\tau]$ 上于每个点 $x_j\in\operatorname{supp}u_s(\cdot,\tau)$ 处满足相容条件 (C$^\pm$)；
\item 若 $\phi$ 无界，则 $u=u_r\in L^\infty(\mathbb{R}\times(\tau,T))$。
\end{enumerate}
\end{theorem}

\begin{remark}\label{remark:ch07-7.3.1}
若 $\phi$ 无界且满足 (A2')，则由文献~\cite{ref67} 定理 1.1 与\cref{theorem:ch07-7.3.3}(ii)，对每个 $u_0\in\mathscr{R}_f^+(\mathbb{R})$ 存在问题 (P) 的唯一熵解，其等待时间 $t_0=0$。事实上，由\cref{theorem:ch07-7.3.3}(ii)，每个所构造熵解 $u$ 是文献~\cite{ref67}意义下的解，因对每个 $\tau\in(0,T)$ 有 $u=u_r\in L^\infty(\mathbb{R}\times(\tau,T))$，且 $\operatorname*{ess\,lim}_{t\to0^+}u(\cdot,t)=u_0$ 在 $\mathscr{R}_f(\mathbb{R})$ 中窄收敛（见\cref{definition:ch05-窄收敛}）。这由\eqref{align:ch07-7.12}与\cref{proposition:ch07-7.1.2}(ii) 得到（另见文献~\cite{ref53} 命题 2，第 38 页）。

当 $\phi$ 有界且 $u_0$ 如 (A0') 时，由文献~\cite{ref17} 命题 3.20，(P) 的每个所构造熵解满足 $u\in C([0,T];\mathscr{R}_f^+(\Omega))$ 且 $u(\cdot,0)=u_0$。故由\cref{theorem:ch07-7.3.3}与\cref{theorem:ch07-唯一性}得下列存在唯一性结果。
\end{remark}

\begin{theorem}\label{theorem:ch07-存在唯一}
设 $\phi$ 有界，(A0')--(A2') 成立。则问题 (P) 存在唯一熵解 $u\in C([0,T];\mathscr{R}_f^+(\Omega))$，它满足相容条件。
\end{theorem}

最后提及下列正则化结果。

\begin{proposition}\label{proposition:ch07-7.3.5}
设 (A0) 与 (A2') 成立。则对几乎所有 $t\in(0,T)$，奇异部分 $u_s(\cdot,t)$ 的支撑集为零测集。
\end{proposition}

\section{存在性结果的证明：逼近问题}\label{section:ch07-逼近问题}

设 (Pn) 的初值 $u_{0n}\in L^1(\mathbb{R})\cap L^\infty(\mathbb{R})$ 满足下列性质（见\cref{lemma:ch06-6.5.1}与\eqref{align:ch6-6.49b}--\eqref{align:ch6-6.49c}）：
\begin{align}
u_{0n}\geqslant0\ \text{在 }\mathbb{R},\qquad
\|u_{0n}\|_{L^1(\mathbb{R})}\leqslant\|u_0\|_{\mathscr{R}_f(\mathbb{R})}, \label{align:ch07-7.44}
\end{align}
\begin{align}
u_{0n}\stackrel{*}{\rightharpoonup}u_0,\quad
u_{0n}\to u_{0r}\ \text{在 }\mathbb{R}\text{ 中几乎处处},\quad
\|u_{0n}-u_{0r}\|_{L^1_{\mathrm{loc}}(\mathbb{R}\setminus\operatorname{supp}u_{0s})}\to0. \label{align:ch07-7.45}
\end{align}
设 $\{u_{0n}^\epsilon\}\subseteq C^\infty_c(\mathbb{R})$，$u_{0n}^\epsilon\geqslant0$，为满足
\begin{align}
\|u_{0n}^\epsilon\|_{L^1(\mathbb{R})}\leqslant\|u_{0n}\|_{L^1(\mathbb{R})}\leqslant\|u_0\|_{\mathscr{R}_f(\mathbb{R})},\quad
\|u_{0n}^\epsilon\|_{L^\infty(\mathbb{R})}\leqslant\|u_{0n}\|_{L^\infty(\mathbb{R})}, \label{align:ch07-7.46}
\end{align}
\begin{align}
u_{0n}^\epsilon\to u_{0n}\ \text{在 }L^1(\mathbb{R}),\quad
u_{0n}^\epsilon\stackrel{*}{\rightharpoonup}u_{0n}\ \text{在 }L^\infty(\mathbb{R})\ \text{当 }\epsilon\to0^+. \label{align:ch07-7.47}
\end{align}
设 $\eta\in C^\infty_c(\mathbb{R})$ 为标准磨光核。置 $\eta_\epsilon(y):=\frac1\epsilon\eta(y/\epsilon)$（$\epsilon>0$），
\begin{align}
\phi_\epsilon(y):=(\eta_\epsilon*\phi)(y)-(\eta_\epsilon*\phi)(0)\quad(y\in\mathbb{R}), \label{align:ch07-7.48}
\end{align}
其中 $\phi:=\phi\chi_{[0,\infty)}$。与 (Pn) 相伴的正则化问题
\begin{align*}
\begin{cases}
\partial_t u_n^\epsilon+\partial_x[\phi_\epsilon(u_n^\epsilon)]=\epsilon\,\partial_{xx}u_n^\epsilon & \text{在 }S,\\
u_n^\epsilon=u_{0n}^\epsilon & \text{在 }\mathbb{R}\times\{0\},
\end{cases}
\qquad(Pn^\epsilon)
\end{align*}
有唯一强解 $u_n^\epsilon\in C([0,T];H^2(\mathbb{R}))\cap L^\infty(S)$，使 $\partial_t u_n^\epsilon\in L^2(S)$ 且 $u_n^\epsilon\geqslant0$（例如见文献~\cite{ref68}）。

族 $\{u_n^\epsilon\}$ 的若干性质收集于下列引理。第一个引理的证明近乎标准（例如见文献~\cite{ref34}），故从略。

\begin{lemma}\label{lemma:ch07-7.4.1}
设 $u_n^\epsilon$ 为问题 (Pn$^\epsilon$) 的解（$\epsilon>0$，$n\in\mathbb{N}$）。则：
\begin{align}
u_n^\epsilon\geqslant0\ \text{在 }S,\qquad
\|u_n^\epsilon\|_{L^\infty(S)}\leqslant\|u_{0n}\|_{L^\infty(\mathbb{R})}, \label{align:ch07-7.49}
\end{align}
\begin{align}
\|u_n^\epsilon(\cdot,t)\|_{L^1(\mathbb{R})}=\|u_{0n}^\epsilon\|_{L^1(\mathbb{R})}\quad(t\in(0,T)), \label{align:ch07-7.50}
\end{align}
\begin{align}
\|u_n^\epsilon\|_{L^\infty(0,T;L^1(\mathbb{R}))}\leqslant\|u_{0n}\|_{L^1(\mathbb{R})}\leqslant\|u_0\|_{\mathscr{R}_f(\mathbb{R})}, \label{align:ch07-7.51}
\end{align}
\begin{align}
\|u_n^\epsilon(\cdot+h,\cdot)-u_n^\epsilon\|_{L^\infty(0,T;L^1(\mathbb{R}))}
\leqslant\|u_{0n}^\epsilon(\cdot+h)-u_{0n}^\epsilon\|_{L^1(\mathbb{R})}\quad\forall h\in\mathbb{R}. \label{align:ch07-7.52}
\end{align}
\end{lemma}

\begin{lemma}\label{lemma:ch07-7.4.2}
设\eqref{align:ch07-7.1}成立。则存在 $C>0$（仅依赖于 $\|u_0\|_{\mathscr{R}_f(\mathbb{R})}$），使对任意 $n\in\mathbb{N}$、$\epsilon\in(0,1)$ 及 $p\in(0,1)$，
\begin{align}
\epsilon\iint_S(1+u_n^\epsilon)^{p-2}(\partial_x u_n^\epsilon)^2\,\mathrm{d}x\,\mathrm{d}t
\leqslant\frac{C}{p(1-p)}. \label{align:ch07-7.53}
\end{align}
\end{lemma}

\begin{proof}
设 $U\in C^2([0,\infty))$，$U'\geqslant0$ 在 $(0,\infty)$。置
\begin{align}
\Lambda_{U,\epsilon}(y):=\int_0^y U'(z)\phi_\epsilon'(z)\,\mathrm{d}z+\vartheta_U\quad(\vartheta_U\in\mathbb{R}). \label{align:ch07-7.54}
\end{align}
由\eqref{align:ch07-7.1}与\eqref{align:ch07-7.48}，对所有 $y\geqslant0$，
\begin{align}
|\Lambda_{U,\epsilon}(y)|\leqslant\int_0^y U'(z)|\phi_\epsilon'(z)|\,\mathrm{d}z+|\vartheta_U|
\leqslant M[U(y)-U(0)]+|\vartheta_U|. \label{align:ch07-7.55}
\end{align}
将 (Pn$^\epsilon$) 第一式乘以 $U'(u_n^\epsilon)$，得
\begin{align}
\partial_t[U(u_n^\epsilon)]+\partial_x[\Lambda_{U,\epsilon}(u_n^\epsilon)]
=\epsilon\,\partial_{xx}[U(u_n^\epsilon)]-\epsilon U''(u_n^\epsilon)(\partial_x u_n^\epsilon)^2
\quad\text{在 }S. \label{align:ch07-7.56}
\end{align}
于是对所有 $\zeta\in C^1([0,T];C^2_c(\mathbb{R}))$，
\begin{align}
\begin{aligned}
\epsilon\iint_S&U''(u_n^\epsilon)(\partial_x u_n^\epsilon)^2\zeta\,\mathrm{d}x\,\mathrm{d}t+\int_{\mathbb{R}}U(u_n^\epsilon(x,T))\zeta(x,T)\,\mathrm{d}x\\
&=\int_{\mathbb{R}}U(u_{0n}^\epsilon)\zeta(x,0)\,\mathrm{d}x
+\iint_S[U(u_n^\epsilon)\partial_t\zeta+\Lambda_{U,\epsilon}(u_n^\epsilon)\partial_x\zeta+\epsilon U(u_n^\epsilon)\partial_{xx}\zeta]\,\mathrm{d}x\,\mathrm{d}t.
\end{aligned} \label{align:ch07-7.57}
\end{align}
取 $\vartheta_U=0$，$U(y)=(1+u)^p-1$（$p\in(0,1)$），及
\begin{align*}
\zeta=\rho_k:=\chi_{[-k,k]}+\rho(\cdot-k)\chi_{[k,k+1)}+\rho(\cdot+k)\chi_{(-k-1,-k]}\quad(k\in\mathbb{N})
\end{align*}
其中 $\rho\in C^2_c((-1,1))$ 满足 $\rho(0)=1$、$0\leqslant\rho\leqslant1$、$\rho'(0)=\rho''(0)=0$。则 $0\leqslant U(y)\leqslant y$（$y\geqslant0$），由\eqref{align:ch07-7.51}、\eqref{align:ch07-7.55}与\eqref{align:ch07-7.57}得
\begin{align*}
\epsilon p(1-p)\iint_S(1+u_n^\epsilon)^{p-2}(\partial_x u_n^\epsilon)^2\rho_k\,\mathrm{d}x\,\mathrm{d}t
\leqslant\int_{\mathbb{R}}u_{0n}^\epsilon(x)\,\mathrm{d}x
+\iint_S[Mu_n^\epsilon|\rho_k'|+\epsilon u_n^\epsilon|\rho_k''|]\,\mathrm{d}x\,\mathrm{d}t
\leqslant[1+(M+1)T\|\rho\|_{C^2([-1,1])}]\|u_0\|_{\mathscr{R}_f(\mathbb{R})}=:C
\end{align*}
对所有 $\epsilon\in(0,1)$ 与 $k\in\mathbb{N}$ 成立。令 $k\to\infty$ 得\eqref{align:ch07-7.53}，结论得证。
\end{proof}

\begin{lemma}\label{lemma:ch07-7.4.3}
设\eqref{align:ch07-7.1}成立，$U\in C^2([0,\infty))$。设存在 $K\geqslant0$ 与 $p\in(0,1)$ 使
\begin{align}
|U''(y)|\leqslant K(1+y)^{p-2}\quad\text{对所有 }y\in[0,\infty). \label{align:ch07-7.58}
\end{align}
则存在 $C_p>0$，使对所有 $n\in\mathbb{N}$ 与 $\epsilon>0$，
\begin{align}
\epsilon\iint_S|U''(u_n^\epsilon)|(\partial_x u_n^\epsilon)^2\,\mathrm{d}x\,\mathrm{d}t\leqslant C_p. \label{align:ch07-7.59}
\end{align}
\end{lemma}

\begin{proof}
由\eqref{align:ch07-7.53}与\eqref{align:ch07-7.58}立即得证。
\end{proof}

\begin{lemma}\label{lemma:ch07-7.4.4}
设\eqref{align:ch07-7.1}成立，$U\in C^2([0,\infty))$，$U'\in L^\infty(0,\infty)$ 满足\eqref{align:ch07-7.58}。对任意 $\rho\in C^2_c(\mathbb{R})$ 与 $t\in(0,T)$，置
\begin{align}
U_{n,\rho}^\epsilon(t):=\int_{\mathbb{R}}U(u_n^\epsilon(x,t))\rho(x)\,\mathrm{d}x,\qquad
U_{n,\rho}(t):=\int_{\mathbb{R}}U(u_n(x,t))\rho(x)\,\mathrm{d}x. \label{align:ch07-7.60}
\end{align}
则族 $\{U_{n,\rho}^\epsilon\}$ 在 $BV(0,T)$ 中有界。
\end{lemma}

\begin{proof}
由等式\eqref{align:ch07-7.56}，
\begin{align}
(U_{n,\rho}^\epsilon)'(t)=\int_{\mathbb{R}}[\Lambda_{U,\epsilon}(u_n^\epsilon)\rho'+\epsilon U(u_n^\epsilon)\rho''-\epsilon U''(u_n^\epsilon)(\partial_x u_n^\epsilon)^2\rho](x,t)\,\mathrm{d}x. \label{align:ch07-7.61}
\end{align}
因 $U'\in L^\infty(0,\infty)$，存在 $N>0$ 使
\begin{align}
|U(y)|\leqslant N(1+y)\quad\text{对所有 }y\geqslant0. \label{align:ch07-7.62}
\end{align}
故由\eqref{align:ch07-7.1}、\eqref{align:ch07-7.48}与\eqref{align:ch07-7.54}，
\begin{align*}
|\Lambda_{U,\epsilon}(u_n^\epsilon)|\leqslant M\|U'\|_\infty|u_n^\epsilon|+|\vartheta_U|=:\widetilde M u_n^\epsilon+|\vartheta_U|,
\end{align*}
由\eqref{align:ch07-7.61}
\begin{align}
\begin{aligned}
|(U_{n,\rho}^\epsilon)'(t)|&\leqslant\|\rho\|_{C^2(\mathbb{R})}\int_{\operatorname{supp}\rho}\{(\widetilde M+\epsilon N)u_n^\epsilon(x,t)+\epsilon N+|\vartheta_U|\}\,\mathrm{d}x\\
&\quad+\epsilon\|\rho\|_{L^\infty(\mathbb{R})}\int_{\mathbb{R}}|U''(u_n^\epsilon)|(\partial_x u_n^\epsilon)^2(x,t)\,\mathrm{d}x.
\end{aligned} \label{align:ch07-7.63}
\end{align}
由\eqref{align:ch07-7.63}、\eqref{align:ch07-7.51}与\eqref{align:ch07-7.59}，存在 $C_{p,\rho}>0$ 使
\begin{align}
\|(U_{n,\rho}^\epsilon)'\|_{L^1(0,T)}\leqslant\|\rho\|_{C^2(\mathbb{R})}[(\widetilde M+N)T\|u_0\|_{\mathscr{R}_f(\mathbb{R})}+C_{p,\rho}]. \label{align:ch07-7.64}
\end{align}
另一方面，由\eqref{align:ch07-7.51}与\eqref{align:ch07-7.62}，
\begin{align}
\|U_{n,\rho}^\epsilon\|_{L^1(0,T)}\leqslant NT\|\rho\|_{L^\infty(\mathbb{R})}[\|u_0\|_{\mathscr{R}_f(\mathbb{R})}+\lambda(\operatorname{supp}\rho)] \label{align:ch07-7.65}
\end{align}
（$\lambda$ 如惯例表示 $\mathbb{R}$ 中 Lebesgue 测度）。由\eqref{align:ch07-7.64}--\eqref{align:ch07-7.65}得结论。
\end{proof}

置 $R_L:=I_L\times(0,T)$，$I_L:=(-L,L)$。由上述引理得下列收敛结果。

\begin{lemma}\label{lemma:ch07-7.4.5}
设 $\phi\in C([0,\infty))$。则存在子列 $\{u_n^{\epsilon_j}\}\subseteq\{u_n^\epsilon\}$（$\epsilon_j\to0^+$，$j\to\infty$）与 $u_n\in L^\infty(S)\cap L^\infty(0,T;L^1(\mathbb{R}))$、$u_n\geqslant0$，使当 $j\to\infty$ 时
\begin{align}
u_n^{\epsilon_j}\stackrel{*}{\rightharpoonup}u_n\quad\text{在 }L^\infty(S), \label{align:ch07-7.66}
\end{align}
\begin{align}
u_n^{\epsilon_j}(\cdot,t)\stackrel{*}{\rightharpoonup}u_n(\cdot,t)\quad\text{在 }L^\infty(\mathbb{R})\ \text{对几乎所有 }t\in(0,T), \label{align:ch07-7.67}
\end{align}
\begin{align}
u_n^{\epsilon_j}(\cdot,t)\to u_n(\cdot,t)\quad\text{在 }L^1(I_L)\ \text{对几乎所有 }t\in(0,T)\text{ 与所有 }L>0, \label{align:ch07-7.68}
\end{align}
\begin{align}
u_n^{\epsilon_j}\to u_n\quad\text{在 }L^1(R_L)\ \text{对所有 }L>0, \label{align:ch07-7.69}
\end{align}
\begin{align}
u_n^{\epsilon_j}\to u_n,\quad \phi_{\epsilon_j}(u_n^{\epsilon_j})\to\phi(u_n)\quad\text{在 }S\text{ 中几乎处处}. \label{align:ch07-7.70}
\end{align}
此外，
\begin{align}
\|u_n\|_{L^\infty(S)}\leqslant\|u_{0n}\|_{L^\infty(\mathbb{R})},\quad
\|u_n\|_{L^\infty(0,T;L^1(\mathbb{R}))}\leqslant\|u_{0n}\|_{L^1(\mathbb{R})}\leqslant\|u_0\|_{\mathscr{R}_f(\mathbb{R})}. \label{align:ch07-7.71}
\end{align}
\end{lemma}

\begin{proof}
\eqref{align:ch07-7.66}中的收敛、$u_n$ 的非负性及\eqref{align:ch07-7.71}中第一个不等式由\eqref{align:ch07-7.49}立即得到。为证\eqref{align:ch07-7.67}，对任意 $\rho\in C^2_c(\mathbb{R})$ 置
\begin{align*}
I_{n,\rho}^{\epsilon_j}(t):=\int_{\mathbb{R}}u_n^{\epsilon_j}(x,t)\rho(x)\,\mathrm{d}x\quad(t\in(0,T)).
\end{align*}
由\cref{lemma:ch07-7.4.4}（取 $U(y)=y$），序列 $\{I_{n,\rho}^{\epsilon_j}\}$ 在 $BV(0,T)$ 中有界。故存在子列（为简便不重新标号）与 $I_{n,\rho}\in BV(0,T)$ 使
\begin{align}
I_{n,\rho}^{\epsilon_j}\to I_{n,\rho}\ \text{在 }L^1(0,T)\text{ 且在 }(0,T)\text{ 中几乎处处}. \label{align:ch07-7.72}
\end{align}
另一方面，由\eqref{align:ch07-7.66}，
\begin{align*}
\int_0^T I_{n,\rho}^{\epsilon_j}(t)\,\mathrm{d}t=\iint_S u_n^{\epsilon_j}(x,t)\rho(x)\,\mathrm{d}x\,\mathrm{d}t
\to\iint_S u_n(x,t)\rho(x)\,\mathrm{d}x\,\mathrm{d}t=\int_0^T\mathrm{d}t\int_{\mathbb{R}}u_n(x,t)\rho(x)\,\mathrm{d}x,
\end{align*}
故 $I_{n,\rho}=\int_{\mathbb{R}}u_n(x,t)\rho(x)\,\mathrm{d}x$（对几乎所有 $t\in(0,T)$），且\eqref{align:ch07-7.72}对整个序列成立。综上，对任意 $\rho\in C^2_c(\mathbb{R})$，存在零测集 $N\subseteq(0,T)$ 使
\begin{align*}
\int_{\mathbb{R}}u_n^{\epsilon_j}(x,t)\rho(x)\,\mathrm{d}x\to\int_{\mathbb{R}}u_n(x,t)\rho(x)\,\mathrm{d}x
\end{align*}
对所有 $t\in(0,T)\setminus N$ 成立。因 $C^2_c(\mathbb{R})$ 可分，$N$ 的选取可与 $\rho$ 无关，故\eqref{align:ch07-7.67}成立。

由\eqref{align:ch07-7.51}、\eqref{align:ch07-7.52}与 Fréchet--Kolmogorov 定理，对所有 $t\in(0,T)$ 与 $L>0$，序列 $\{u_n^{\epsilon_j}(\cdot,t)\}$ 在 $L^1(I_L)$ 中相对紧。如前论证，由\eqref{align:ch07-7.67}即得\eqref{align:ch07-7.68}。由\eqref{align:ch07-7.51}与\eqref{align:ch07-7.68}得
\begin{align*}
\|u_n\|_{L^\infty(0,T;L^1(I_L))}\leqslant\|u_{0n}\|_{L^1(\mathbb{R})}\leqslant\|u_0\|_{\mathscr{R}_f(\mathbb{R})}\quad\forall L>0,
\end{align*}
由 $L$ 任意得\eqref{align:ch07-7.71}中第二个不等式。由\eqref{align:ch07-7.51}与\eqref{align:ch07-7.68}及控制收敛定理得\eqref{align:ch07-7.69}中收敛。最后，\eqref{align:ch07-7.70}中第一个收敛由\eqref{align:ch07-7.69}（可能抽取子列，不重新标号）得到。因对任意紧集 $K\subseteq\mathbb{R}$，当 $\epsilon\to0^+$ 时 $\|\phi_\epsilon-\phi\|_{C(K)}\to0$，\eqref{align:ch07-7.70}中第二个收敛亦得证。
\end{proof}

\begin{lemma}\label{lemma:ch07-7.4.6}
设\eqref{align:ch07-7.1}成立，$U\in C^2([0,\infty))$，$U'\in L^\infty(0,\infty)$ 满足\eqref{align:ch07-7.58}。对任意 $\rho\in C^2_c(\mathbb{R})$ 与 $t\in(0,T)$，置
\begin{align}
U_{n,\rho}(t):=\int_{\mathbb{R}}U(u_n(x,t))\rho(x)\,\mathrm{d}x \label{align:ch07-7.73}
\end{align}
其中 $u_n\in L^\infty(S)\cap L^\infty(0,T;L^1(\mathbb{R}))$ 由\cref{lemma:ch07-7.4.5}给出。则：
\begin{enumerate}[(i)]
\item 存在子列 $\{U_{n,\rho}^{\epsilon_j}\}\subseteq\{U_{n,\rho}^\epsilon\}$（$\epsilon_j\to0^+$）（$U_{n,\rho}^\epsilon$ 由\eqref{align:ch07-7.60}定义）使
\begin{align}
U_{n,\rho}^{\epsilon_j}\to U_{n,\rho}\quad\text{在 }L^1(0,T); \label{align:ch07-7.74}
\end{align}
\item 族 $\{U_{n,\rho}\}$ 在 $BV(0,T)$ 中有界。
\end{enumerate}
\end{lemma}

\begin{proof}
关于 (i)，设 $\{u_n^{\epsilon_j}\}$ 为\cref{lemma:ch07-7.4.5}中子列。则对几乎所有 $t\in(0,T)$，由\eqref{align:ch07-7.70}中第一个收敛有 $U(u_n^{\epsilon_j}(\cdot,t))\to U(u_n(\cdot,t))$ 在 $\mathbb{R}$ 中几乎处处，由控制收敛定理得\eqref{align:ch07-7.74}。关于 (ii)，由\eqref{align:ch07-7.74}与\eqref{align:ch07-7.65}得
\begin{align}
\|U_{n,\rho}\|_{L^1(0,T)}=\lim_{j\to\infty}\|U_{n,\rho}^{\epsilon_j}\|_{L^1(0,T)}
\leqslant NT\|\rho\|_{L^\infty(\mathbb{R})}[\|u_0\|_{\mathscr{R}_f(\mathbb{R})}+\lambda(\operatorname{supp}\rho)]. \label{align:ch07-7.75}
\end{align}
另一方面，由\eqref{align:ch07-7.64}与 $L^1(0,T)$ 拓扑下总变差的下半连续性（见\cref{remark:ch03-3.1.4}(ii)），
\begin{align}
\|U_{n,\rho}'\|_{\mathscr{R}_f(0,T)}\leqslant\|\rho\|_{C^2(\mathbb{R})}[(\widetilde M+N)T\|u_0\|_{\mathscr{R}_f(\mathbb{R})}+C_{p,\rho}]. \label{align:ch07-7.76}
\end{align}
由\eqref{align:ch07-7.75}--\eqref{align:ch07-7.76}得 (ii)。
\end{proof}

回忆下列定义（例如见文献~\cite{ref34}）。

\begin{definition}\label{definition:ch07-Pn熵解}
函数 $u_n\in L^\infty(0,T;L^1(\mathbb{R}))\cap L^\infty(S)$ 称为问题 (Pn) 的熵解，若对每个 $\zeta\in C^1([0,T];C^1_c(\mathbb{R}))$、$\zeta(\cdot,T)=0$ 在 $\mathbb{R}$ 上、$\zeta\geqslant0$ 及任意 $l\in[0,\infty)$，有
\begin{align}
\iint_S[E_l(u_n)\partial_t\zeta+F_l(u_n)\partial_x\zeta]\,\mathrm{d}x\,\mathrm{d}t
+\int_{\mathbb{R}}E_l(u_{0n})\zeta(x,0)\,\mathrm{d}x\geqslant0, \label{align:ch07-7.77}
\end{align}
其中 $E_l(y)=|y-l|$，$F_l(y)=\operatorname{sgn}(y-l)[\phi(y)-\phi(l)]$（$y\in[0,\infty)$）。
\end{definition}

\begin{remark}\label{remark:ch07-7.4.1}
熵解是弱解：若 $\zeta\in C^1([0,T];C^1_c(\mathbb{R}))$、$\zeta(\cdot,T)=0$ 在 $\mathbb{R}$ 上，则
\begin{align}
\iint_S[u_n\partial_t\zeta+\phi(u_n)\partial_x\zeta]\,\mathrm{d}x\,\mathrm{d}t
+\int_{\mathbb{R}}u_{0n}\zeta(x,0)\,\mathrm{d}x=0. \label{align:ch07-7.78}
\end{align}
\end{remark}

\begin{proposition}\label{proposition:ch07-7.4.7}
设 $\phi\in C([0,\infty))$。则\cref{lemma:ch07-7.4.5}给出的函数 $u_n$ 是问题 (Pn)（$n\in\mathbb{N}$）的熵解。此外，
\begin{align}
\|u_n(\cdot,t)\|_{L^1(\mathbb{R})}=\|u_{0n}\|_{L^1(\mathbb{R})}\quad\text{对几乎所有 }t\in(0,T) \label{align:ch07-7.79}
\end{align}
且
\begin{align}
\|u_n(\cdot+h,\cdot)-u_n\|_{L^\infty(0,T;L^1(\mathbb{R}))}
\leqslant\|u_{0n}(\cdot+h)-u_{0n}\|_{L^1(\mathbb{R})}\quad\forall h\in\mathbb{R}. \label{align:ch07-7.80}
\end{align}
若 $\phi$ 局部 Lipschitz 连续，则 $u_n$ 是 (Pn) 的唯一熵解。
\end{proposition}

\begin{proof}
设 $\zeta$ 如\cref{definition:ch07-Pn熵解}，$E_l(y)=|y-l|$，$F_{l,\epsilon}(y)=\operatorname{sgn}(y-l)[\phi_\epsilon(y)-\phi_\epsilon(l)]$（$y,l\in[0,\infty)$）。由 (Pn$^\epsilon$) 作标准计算得
\begin{align}
\iint_S\{E_l(u_n^\epsilon)(\partial_t\zeta+\epsilon\zeta_{xx})+F_{l,\epsilon}(u_n^\epsilon)\partial_x\zeta\}\,\mathrm{d}x\,\mathrm{d}t
+\int_{\mathbb{R}}E_l(u_{0n}^\epsilon)\zeta(x,0)\,\mathrm{d}x\geqslant0. \label{align:ch07-7.81}
\end{align}
设 $\{u_n^{\epsilon_j}\}\subseteq\{u_n^\epsilon\}$ 由\cref{lemma:ch07-7.4.5}给出。由\eqref{align:ch07-7.48}--\eqref{align:ch07-7.49}，
\begin{align*}
\|\phi_{\epsilon_j}(u_n^{\epsilon_j})\|_{L^\infty(S)}
\leqslant\sup_{|y|\leqslant\|u_{0n}\|_{L^\infty(\mathbb{R})}}|\phi_{\epsilon_j}(y)|
\leqslant\sup_{|y|\leqslant\|u_{0n}\|_{L^\infty(\mathbb{R})}}|\phi(y)|.
\end{align*}
由\eqref{align:ch07-7.70}、上式与控制收敛定理，
\begin{align*}
\iint_S F_{l,\epsilon_j}(u_n^{\epsilon_j})\partial_x\zeta\,\mathrm{d}x\,\mathrm{d}t
\to\iint_S F_l(u_n)\partial_x\zeta\,\mathrm{d}x\,\mathrm{d}t.
\end{align*}
由\eqref{align:ch07-7.69}与\eqref{align:ch07-7.47}亦有
\begin{align*}
\iint_S E_l(u_n^{\epsilon_j})(\partial_t\zeta+\epsilon_j\zeta_{xx})\,\mathrm{d}x\,\mathrm{d}t
\to\iint_S E_l(u_n)\partial_t\zeta\,\mathrm{d}x\,\mathrm{d}t,
\end{align*}
\begin{align*}
\int_{\mathbb{R}}E_l(u_{0n}^{\epsilon_j})\zeta(x,0)\,\mathrm{d}x\to\int_{\mathbb{R}}E_l(u_{0n})\zeta(x,0)\,\mathrm{d}x.
\end{align*}
在\eqref{align:ch07-7.81}（取 $\epsilon=\epsilon_j$）中令 $j\to\infty$ 得\eqref{align:ch07-7.77}，故 $u_n$ 是 (Pn) 的熵解。唯一性由 Kružkov 唯一性定理得到（例如见文献~\cite{ref91}）。

不等式\eqref{align:ch07-7.80}由\eqref{align:ch07-7.52}与\eqref{align:ch07-7.68}得到。关于\eqref{align:ch07-7.79}，由\eqref{align:ch07-7.78}经标准论证得，对几乎所有 $t\in(0,T)$ 与所有 $\rho\in C^1_c(\mathbb{R})$，
\begin{align}
\int_{\mathbb{R}}u_n(x,t)\rho(x)\,\mathrm{d}x-\int_{\mathbb{R}}u_{0n}(x)\rho(x)\,\mathrm{d}x
=\int_0^t\int_{\mathbb{R}}\phi(u_n)(x,s)\rho'(x)\,\mathrm{d}x\,\mathrm{d}s. \label{align:ch07-7.82}
\end{align}
设 $\rho_k\in C^1_c(\mathbb{R})$（$k\in\mathbb{N}$）满足
\begin{align*}
\rho_k(x)=\begin{cases}1&\text{若 }|x|\leqslant k,\\0&\text{若 }|x|\geqslant k+1,\end{cases}
\end{align*}
且 $\|\rho_k'\|_{L^\infty(\mathbb{R})}\leqslant2$。因 $u_n\in L^1(S)$，在\eqref{align:ch07-7.82}中取 $\rho=\rho_k$ 并令 $k\to\infty$，利用\eqref{align:ch07-7.1}得
\begin{align*}
\left|\int_0^t\int_{\mathbb{R}}\phi(u_n)(x,s)\rho_k'(x)\,\mathrm{d}x\,\mathrm{d}s\right|
\leqslant2M\int_0^t\int_{\{k\leqslant|x|\leqslant k+1\}}|u_n(x,s)|\,\mathrm{d}x\,\mathrm{d}s\to0.
\end{align*}
另一方面，由单调收敛定理，
\begin{align*}
\int_{\mathbb{R}}u_n(x,t)\rho_k(x)\,\mathrm{d}x\to\int_{\mathbb{R}}u_n(x,t)\,\mathrm{d}x,\quad
\int_{\mathbb{R}}u_{0n}(x)\rho_k(x)\,\mathrm{d}x\to\int_{\mathbb{R}}u_{0n}(x)\,\mathrm{d}x.
\end{align*}
由\eqref{align:ch07-7.82}得等式\eqref{align:ch07-7.79}。
\end{proof}

为后文引用，证明下列结果。

\begin{lemma}\label{lemma:ch07-7.4.8}
设\eqref{align:ch07-7.1}成立，$U\in C^2([0,\infty))\cap L^\infty(0,\infty)$ 满足\eqref{align:ch07-7.58}。置
\begin{align}
\Lambda_U(y):=\int_0^y U'(z)\phi'(z)\,\mathrm{d}z+\vartheta_U\quad(y,\vartheta_U\in\mathbb{R}), \label{align:ch07-7.83}
\end{align}
并设 $\Lambda_U\in L^\infty(0,\infty)$。设 $u_n$ 为由\cref{proposition:ch07-7.4.7}给出的问题 (Pn)（$n\in\mathbb{N}$）的熵解。则对任意 $n\in\mathbb{N}$，存在 $\sigma_n\in\mathscr{R}_f(S)$，使对所有 $\zeta\in C^1_c(S)$，
\begin{align}
\iint_S\{U(u_n)\partial_t\zeta+\Lambda_U(u_n)\partial_x\zeta\}\,\mathrm{d}x\,\mathrm{d}t=\langle\sigma_n,\zeta\rangle. \label{align:ch07-7.84}
\end{align}
此外有
\begin{align}
\|\sigma_n\|_{\mathscr{R}_f(S)}\leqslant C_p \label{align:ch07-7.85}
\end{align}
其中 $C_p>0$ 如\eqref{align:ch07-7.59}。
\end{lemma}

\begin{proof}
由不等式\eqref{align:ch07-7.59}，对任意固定的 $n\in\mathbb{N}$，族 $\{\epsilon U''(u_n^\epsilon)(\partial_x u_n^\epsilon)^2\}$ 在 $L^1(S)\subseteq\mathscr{R}_f(S)$ 中有界。故由 Banach--Alaoglu 定理，存在序列 $\{\epsilon_k\}$（$\epsilon_k\to0^+$，$k\to\infty$）与 $\sigma_n\in\mathscr{R}_f(S)$ 使
\begin{align}
\epsilon_k U''(u_n^{\epsilon_k})(\partial_x u_n^{\epsilon_k})^2\stackrel{*}{\rightharpoonup}\sigma_n\quad\text{当 }k\to\infty. \label{align:ch07-7.86}
\end{align}
由范数的下半连续性（见\cref{remark:ch05-5.1.3}），从\eqref{align:ch07-7.86}得不等式\eqref{align:ch07-7.85}。

设 $\zeta\in C^2_c(S)$。由\eqref{align:ch07-7.56}，
\begin{align}
\epsilon\iint_S U''(u_n^\epsilon)(\partial_x u_n^\epsilon)^2\zeta\,\mathrm{d}x\,\mathrm{d}t
=\iint_S[U(u_n^\epsilon)\partial_t\zeta+\Lambda_{U,\epsilon}(u_n^\epsilon)\partial_x\zeta+\epsilon U(u_n^\epsilon)\partial_{xx}\zeta]\,\mathrm{d}x\,\mathrm{d}t \label{align:ch07-7.87}
\end{align}
其中 $\Lambda_{U,\epsilon}$ 由\eqref{align:ch07-7.54}定义。由\eqref{align:ch07-7.1}与\eqref{align:ch07-7.49}中第二个不等式，对任意固定的 $n\in\mathbb{N}$，
\begin{align*}
|\Lambda_{U,\epsilon}(u_n^\epsilon)|\leqslant M\int_0^{\|u_{0n}\|_{L^\infty(\mathbb{R})}}|U'(z)|\,\mathrm{d}z+|\vartheta_U|<\infty,
\end{align*}
故对固定的 $n$，族 $\{\Lambda_{U,\epsilon}(u_n^\epsilon)\}$ 在 $L^\infty(S)$ 中有界。在\eqref{align:ch07-7.87}（取 $\epsilon=\epsilon_k$）中令 $k\to\infty$，如\cref{proposition:ch07-7.4.7}证明中那样论证，即得等式\eqref{align:ch07-7.84}。
\end{proof}

\section{存在性结果的证明}\label{section:ch07-存在性证明}

\cref{theorem:ch07-7.3.1}的证明依赖下列结果，它是\cref{theorem:ch06-6.6.1}的对应物。如\cref{theorem:ch06-6.6.1}情形，证明需要若干预备结果。

\begin{theorem}\label{theorem:ch07-7.5.1}
设 (A0)--(A1) 成立，$\{u_n\}$ 为问题 (Pn) 的熵解序列。则存在子列 $\{u_{n_k}\}\subseteq\{u_n\}$，具有下列性质：
\begin{enumerate}[(i)]
\item 存在 $u\in L^\infty_{w^*}(0,T;\mathscr{R}_f^+(\mathbb{R}))$ 使
\begin{align}
u_{n_k}\stackrel{*}{\rightharpoonup}u\quad\text{在 }L^\infty_{w^*}(0,T;\mathscr{R}_f(\mathbb{R})); \label{align:ch07-7.88}
\end{align}
\item 与 $\{u_{n_k}\}$ 相伴的 Young 测度序列 $\{\nu_{n_k}\}\subseteq\mathcal{Y}^+(S;\mathbb{R})$ 窄收敛于 Young 测度 $\nu\in\mathcal{Y}^+(S;\mathbb{R})$，且对几乎所有 $(x,t)\in S$，
\begin{align}
u_r(x,t)=\int_{[0,\infty)}y\,\mathrm{d}\nu_{(x,t)}(y), \label{align:ch07-7.89}
\end{align}
\begin{align}
\phi(u_r(x,t))=\int_{[0,\infty)}\phi(y)\,\mathrm{d}\nu_{(x,t)}(y), \label{align:ch07-7.90}
\end{align}
其中 $\{\nu_{(x,t)}\}_{(x,t)\in S}$ 为 $\nu$ 的离散分解；
\item 对每个满足
\begin{align}
\lim_{y\to\infty}\frac{U(y)}{y}=:M_U\in\mathbb{R} \label{align:ch07-7.91}
\end{align}
的 $U\in C([0,\infty))$ 与每个 $\zeta\in C([0,T];C_c(\mathbb{R}))$，有
\begin{align}
\lim_{k\to\infty}\iint_S U(u_{n_k})\zeta\,\mathrm{d}x\,\mathrm{d}t
=\iint_S U^*\zeta\,\mathrm{d}x\,\mathrm{d}t+M_U\int_0^T\langle u_s(\cdot,t),\zeta(\cdot,t)\rangle\,\mathrm{d}t, \label{align:ch07-7.92}
\end{align}
其中 $U^*\in L^\infty(0,T;L^1_{\mathrm{loc}}(\mathbb{R}))$ 定义为
\begin{align}
U^*(x,t):=\int_{[0,\infty)}U(y)\,\mathrm{d}\nu_{(x,t)}(y)
\quad\text{对几乎所有 }(x,t)\in S. \label{align:ch07-7.93}
\end{align}
\end{enumerate}
\end{theorem}

为证\cref{theorem:ch07-7.5.1}，需要下列结果。

\begin{proposition}\label{proposition:ch07-7.5.2}
设 (A0)--(A1) 成立，$\{u_n\}$ 为问题 (Pn) 的熵解序列。则存在子列 $\{u_{n_k}\}\subseteq\{u_n\}$、Young 测度 $\nu\in\mathcal{Y}^+(S;\mathbb{R})$ 与 $\sigma\in L^\infty_{w^*}(0,T;\mathscr{R}_f^+(\mathbb{R}))$，使
\begin{align}
u_{n_k}\stackrel{*}{\rightharpoonup}u_b+\sigma\quad\text{在 }L^\infty_{w^*}(0,T;\mathscr{R}_f(\mathbb{R})), \label{align:ch07-7.94}
\end{align}
其中 $u_b\in L^\infty(0,T;L^1(\mathbb{R}))$ 定义为
\begin{align}
u_b(x,t):=\int_{[0,\infty)}y\,\mathrm{d}\nu_{(x,t)}(y)
\quad\text{对几乎所有 }(x,t)\in S. \label{align:ch07-7.95}
\end{align}
\end{proposition}

\begin{proof}
由\eqref{align:ch07-7.71}中第二个不等式，序列 $\{u_n\}$ 在 $L^\infty(0,T;L^1(\mathbb{R}))\subseteq L^\infty_{w^*}(0,T;\mathscr{R}_f(\mathbb{R}))$ 中有界。由\cref{proposition:ch04-4.4.16}与 Banach--Alaoglu 定理，存在 $\mu\in L^\infty_{w^*}(0,T;\mathscr{R}_f^+(\mathbb{R}))$ 使
\begin{align}
u_{n_k}\stackrel{*}{\rightharpoonup}\mu\quad\text{在 }L^\infty_{w^*}(0,T;\mathscr{R}_f(\mathbb{R})). \label{align:ch07-7.96}
\end{align}
另一方面，因 $\{u_{n_k}\}$ 在 $L^1(S)$ 中有界，由\cref{remark:ch05-5.4.4}(ii)，存在 $\{u_{n_k}\}$ 的子列（不重新标号）与 $\sigma\in\mathscr{R}_f^+(S)$ 使
\begin{align}
u_{n_k}\stackrel{*}{\rightharpoonup}u_b+\sigma\quad\text{在 }\mathscr{R}_f(S). \label{align:ch07-7.97}
\end{align}
由\eqref{align:ch07-7.96}--\eqref{align:ch07-7.97}得 $\mu=u_b+\sigma$，结论得证。
\end{proof}

\begin{proposition}\label{proposition:ch07-7.5.3}
设 (A0)--(A1) 成立，$\{u_n\}$ 为问题 (Pn) 的熵解序列。设 $\{u_{n_k}\}\subseteq\{u_n\}$、$\nu\in\mathcal{Y}^+(S;\mathbb{R})$ 与 $\sigma\in L^\infty_{w^*}(0,T;\mathscr{R}_f^+(\mathbb{R}))$ 由\cref{proposition:ch07-7.5.2}给出。设 $U\in C([0,\infty))$ 满足\eqref{align:ch07-7.91}。则对任意 $\zeta\in C([0,T];C_c(\mathbb{R}))$，
\begin{align}
\lim_{k\to\infty}\iint_S U(u_{n_k})\zeta\,\mathrm{d}x\,\mathrm{d}t
=\iint_S U^*\zeta\,\mathrm{d}x\,\mathrm{d}t+M_U\int_0^T\langle\sigma(\cdot,t),\zeta(\cdot,t)\rangle\,\mathrm{d}t \label{align:ch07-7.98}
\end{align}
其中 $U^*\in L^\infty(0,T;L^1_{\mathrm{loc}}(\mathbb{R}))$ 由\eqref{align:ch07-7.93}给出。
\end{proposition}

\begin{proof}
固定任意 $L>0$。由\eqref{align:ch07-7.91}，对任意 $\epsilon>0$，存在 $y_\epsilon>0$ 使对所有 $y>y_\epsilon$，
\begin{align}
-\epsilon y<U(y)-M_U y<\epsilon y. \label{align:ch07-7.99}
\end{align}
对任意 $j\in\mathbb{N}$、$j\geqslant[y_\epsilon]+1$，设 $f_{j1},f_{j2}\in C([0,\infty))$ 满足
\begin{align}
\begin{cases}
0\leqslant f_{j1}\leqslant1,\ 0\leqslant f_{j2}\leqslant1,\ f_{j1}+f_{j2}=1\text{ 在 }[0,\infty),\\
\operatorname{supp}f_{j1}\subseteq[0,j+1],\ \operatorname{supp}f_{j2}\subseteq[j,\infty).
\end{cases} \label{align:ch07-7.100}
\end{align}
利用\eqref{align:ch07-7.99}与\eqref{align:ch07-7.71}，对固定的 $k\in\mathbb{N}$ 与 $t\in(0,T)$，
\begin{align*}
\begin{aligned}
\int_{I_L}|U(u_{n_k}(x,t))|\,\mathrm{d}x
&=\int_{I_L}f_{j1}(u_{n_k}(x,t))|U(u_{n_k}(x,t))|\,\mathrm{d}x+\int_{I_L}f_{j2}(u_{n_k}(x,t))|U(u_{n_k}(x,t))|\,\mathrm{d}x\\
&\leqslant2L\|U\|_{L^\infty([0,j+1])}+\int_{\{u_{n_k}(\cdot,t)>y_\epsilon\}}|U(u_{n_k}(x,t))|\,\mathrm{d}x\\
&\leqslant2L\|U\|_{L^\infty([0,j+1])}+(|M_U|+\epsilon)\int_{I_L}|u_{n_k}(x,t)|\,\mathrm{d}x\\
&\leqslant2L\|U\|_{L^\infty([0,j+1])}+(|M_U|+\epsilon)\|u_0\|_{\mathscr{R}_f(\mathbb{R})}
\end{aligned}
\end{align*}
（回忆 $I_L=(-L,L)$），故
\begin{align}
\|U(u_{n_k})\|_{L^\infty(0,T;L^1(I_L))}\leqslant2L\|U\|_{L^\infty([0,j+1])}+(|M_U|+\epsilon)\|u_0\|_{\mathscr{R}_f(\mathbb{R})}. \label{align:ch07-7.101}
\end{align}
由 Banach--Alaoglu 定理，存在 $\{U(u_{n_k})\}$ 的子列（不重新标号）与 $\mu\in L^\infty_{w^*}(0,T;\mathscr{R}_f(I_L))$ 使
\begin{align}
U(u_{n_k})\stackrel{*}{\rightharpoonup}\mu\quad\text{在 }L^\infty_{w^*}(0,T;\mathscr{R}_f(I_L)). \label{align:ch07-7.102}
\end{align}
此外，如\cref{proposition:ch05-5.4.13}证明中那样论证，
\begin{align}
U(u_{n_k})\stackrel{*}{\rightharpoonup}U^*+M_U\sigma\quad\text{在 }\mathscr{R}_f(R_L). \label{align:ch07-7.103}
\end{align}
由\eqref{align:ch07-7.102}--\eqref{align:ch07-7.103}，对任意 $\zeta\in C([0,T];C_c(I_L))$，
\begin{align}
\lim_{k\to\infty}\iint_{R_L}U(u_{n_k})\zeta\,\mathrm{d}x\,\mathrm{d}t
=\iint_{R_L}U^*\zeta\,\mathrm{d}x\,\mathrm{d}t+M_U\int_0^T\langle\sigma(\cdot,t),\zeta(\cdot,t)\rangle\,\mathrm{d}t \label{align:ch07-7.104}
\end{align}
其中 $U^*$ 由\eqref{align:ch07-7.93}在 $R_L$ 中几乎处处定义。

等式\eqref{align:ch07-7.98}由\eqref{align:ch07-7.104}得到：因对任意 $\zeta\in C([0,T];C_c(\mathbb{R}))$，可取 $L>0$ 充分大使 $\operatorname{supp}\zeta\subseteq((-L,L)\times[0,T])$。此外，由\eqref{align:ch07-7.91}存在 $M>0$ 使 $|U(y)|\leqslant M(y+1)$（$y\geqslant0$），故 $|U^*|\leqslant M(u_b+1)$ 在 $S$ 中几乎处处（见\eqref{align:ch07-7.93}与\eqref{align:ch07-7.95}）。因 $u_b\in L^\infty(0,T;L^1(\mathbb{R}))$，得 $U^*\in L^\infty(0,T;L^1_{\mathrm{loc}}(\mathbb{R}))$。
\end{proof}

为继续，需要下列引理，其证明在本节末给出。

\begin{lemma}\label{lemma:ch07-7.5.4}
设 (A0)--(A1) 成立，$U\in C([0,\infty))$ 满足\eqref{align:ch07-7.91}。设 $\{u_{n_k}\}$、$\sigma$ 与 $U^*\in L^\infty(0,T;L^1_{\mathrm{loc}}(\mathbb{R}))$ 如\cref{proposition:ch07-7.5.3}。则对每个 $\zeta\in C([0,T];C_c(\mathbb{R}))$，
\begin{align}
\lim_{k\to\infty}\int_0^T\left|\int_{\mathbb{R}}U(u_{n_k}(x,t))\zeta(x,t)\,\mathrm{d}x
-\int_{\mathbb{R}}U^*(x,t)\zeta(x,t)\,\mathrm{d}x
-M_U\langle\sigma(\cdot,t),\zeta(\cdot,t)\rangle\right|\mathrm{d}t=0. \label{align:ch07-7.105}
\end{align}
\end{lemma}

\begin{remark}\label{remark:ch07-7.5.1}
取 $\zeta(x,t)=\rho(x)$、$\rho\in C_c(\mathbb{R})$，由\eqref{align:ch07-7.105}得
\begin{align}
\lim_{k\to\infty}\int_0^T\left|\int_{\mathbb{R}}U(u_{n_k}(x,t))\rho(x)\,\mathrm{d}x-U_\rho^*(t)-M_U\langle\sigma(\cdot,t),\rho\rangle\right|\mathrm{d}t=0, \label{align:ch07-7.106a}
\end{align}
其中
\begin{align}
U_\rho^*(t):=\int_{\mathbb{R}}U^*(x,t)\rho(x)\,\mathrm{d}x. \label{align:ch07-7.106b}
\end{align}
由\eqref{align:ch07-7.106a}与 $C_c(\mathbb{R})$ 的可分性，对每个 $L>0$ 与如\cref{lemma:ch07-7.5.4}的 $U$，存在零测集 $N\subseteq(0,T)$ 与 $\{u_{n_k}\}$ 的子列（不重新标号），使对所有 $t\in(0,T)\setminus N$，
\begin{align}
U(u_{n_k})(\cdot,t)\stackrel{*}{\rightharpoonup}U^*(\cdot,t)+M_U\sigma(\cdot,t)\quad\text{在 }\mathscr{R}_f(I_L). \label{align:ch07-7.107}
\end{align}
\end{remark}

\begin{proposition}\label{proposition:ch07-7.5.5}
设 (A0)--(A1) 成立，$\{u_n\}$ 为问题 (Pn) 的熵解序列。设 $\{u_{n_k}\}\subseteq\{u_n\}$ 与 $\sigma\in L^\infty_{w^*}(0,T;\mathscr{R}_f^+(\mathbb{R}))$ 由\cref{proposition:ch07-7.5.2}给出，并按
\begin{align}
u:=u_b+\sigma \label{align:ch07-7.108}
\end{align}
定义 $u\in L^\infty_{w^*}(0,T;\mathscr{R}_f^+(\mathbb{R}))$。则
\begin{align}
u_r=u_b\quad\text{在 }S\text{ 中几乎处处}, \label{align:ch07-7.109a}
\end{align}
\begin{align}
u_s=\sigma\quad\text{在 }L^\infty_{w^*}(0,T;\mathscr{R}_f(\mathbb{R})). \label{align:ch07-7.109b}
\end{align}
\end{proposition}

\begin{proof}
对任意非负 $\zeta\in C^1([0,T];C^1_c(\mathbb{R}))$、$\zeta(\cdot,T)=0$，将\eqref{align:ch07-7.77}与\eqref{align:ch07-7.78}相加，得
\begin{align*}
\iint_S\{U_j(u_{n_k})\partial_t\zeta+\Lambda_{U_j}(u_{n_k})\partial_x\zeta\}\,\mathrm{d}x\,\mathrm{d}t
\geqslant-\int_{\mathbb{R}}U_j(u_{0n_k})(x)\zeta(x,0)\,\mathrm{d}x,
\end{align*}
其中对所有 $j\in\mathbb{N}$，
\begin{align}
U_j(y):=(y-j)\chi_{[j,\infty)}(y),\quad
\Lambda_{U_j}(y):=\int_0^y U_j'(z)\phi'(z)\,\mathrm{d}z=[\phi(y)-\phi(j)]\chi_{[j,\infty)}(y). \label{align:ch07-7.110}
\end{align}
由标准逼近论证，对上述每个 $\zeta$ 与几乎所有 $\tau\in(0,T)$，
\begin{align}
\begin{aligned}
\int_{\mathbb{R}}&U_j(u_{n_k})(x,\tau)\zeta(x,\tau)\,\mathrm{d}x-\int_{\mathbb{R}}U_j(u_{0n_k})(x)\zeta(x,0)\,\mathrm{d}x\\
&\leqslant\iint_{S_\tau}\{U_j(u_{n_k})\partial_t\zeta+\Lambda_{U_j}(u_{n_k})\partial_x\zeta\}\,\mathrm{d}x\,\mathrm{d}t,
\end{aligned} \label{align:ch07-7.111}
\end{align}
其中 $S_\tau:=\mathbb{R}\times(0,\tau)$。因 $M_{U_j}=1$、$M_{\Lambda_{U_j}}=0$，由\eqref{align:ch07-7.105}得
\begin{align}
\lim_{k\to\infty}\int_0^\tau\left|\int_{\mathbb{R}}[U_j(u_{n_k})\partial_t\zeta](x,t)\,\mathrm{d}x
-\int_{\mathbb{R}}[U_j^*\partial_t\zeta](x,t)\,\mathrm{d}x
-\langle\sigma(\cdot,t),\partial_t\zeta(\cdot,t)\rangle\right|\mathrm{d}t=0, \label{align:ch07-7.112a}
\end{align}
\begin{align}
\lim_{k\to\infty}\int_0^\tau\left|\int_{\mathbb{R}}[\Lambda_{U_j}(u_{n_k})\partial_x\zeta](x,t)\,\mathrm{d}x
-\int_{\mathbb{R}}[\Lambda_{U_j}^*\partial_x\zeta](x,t)\,\mathrm{d}x\right|\mathrm{d}t=0, \label{align:ch07-7.112b}
\end{align}
其中
\begin{align*}
U_j^*:=\int_{[0,\infty)}U_j(y)\,\mathrm{d}\nu(y),\qquad
\Lambda_{U_j}^*:=\int_{[0,\infty)}\Lambda_{U_j}(y)\,\mathrm{d}\nu(y).
\end{align*}
由\eqref{align:ch07-7.112a}得
\begin{align}
\begin{aligned}
\lim_{k\to\infty}\iint_{S_\tau}&\{U_j(u_{n_k})\partial_t\zeta+\Lambda_{U_j}(u_{n_k})\partial_x\zeta\}\,\mathrm{d}x\,\mathrm{d}t\\
&=\iint_{S_\tau}\{U_j^*\partial_t\zeta+\Lambda_{U_j}^*\partial_x\zeta\}\,\mathrm{d}x\,\mathrm{d}t
+\int_0^\tau\langle\sigma(\cdot,t),\partial_t\zeta(\cdot,t)\rangle\,\mathrm{d}t.
\end{aligned} \label{align:ch07-7.113}
\end{align}
另一方面，由\eqref{align:ch07-7.107}与对角论证，存在零测集 $N\subseteq(0,T)$ 与 $\{u_{n_k}\}$ 的子列（不重新标号），使对所有 $\tau\in(0,T)\setminus N$ 与 $j\in\mathbb{N}$，
\begin{align}
\lim_{k\to\infty}\int_{\mathbb{R}}U_j(u_{n_k})(x,\tau)\zeta(x,\tau)\,\mathrm{d}x
=\int_{\mathbb{R}}U_j^*(x,\tau)\zeta(x,\tau)\,\mathrm{d}x+\langle\sigma(\cdot,\tau),\zeta(x,\tau)\rangle. \label{align:ch07-7.114}
\end{align}
此外，因序列 $\{U_j(u_{0n_k})-u_{0n_k}\}$ 在 $L^\infty(\mathbb{R})$ 中有界且
\begin{align*}
\lim_{k\to\infty}(U_j(u_{0n_k})-u_{0n_k})=U_j(u_{0r})-u_{0r}\ \text{在 }\mathbb{R}\text{ 中几乎处处}
\end{align*}
（见\eqref{align:ch07-7.45}），由\eqref{align:ch07-7.45}得
\begin{align}
\lim_{k\to\infty}\int_{\mathbb{R}}U_j(u_{0n_k})(x)\zeta(x,0)\,\mathrm{d}x
=\int_{\mathbb{R}}U_j(u_{0r})(x)\zeta(x,0)\,\mathrm{d}x+\langle u_{0s},\zeta(\cdot,0)\rangle. \label{align:ch07-7.115}
\end{align}
在\eqref{align:ch07-7.111}中令 $k\to\infty$，由\eqref{align:ch07-7.113}--\eqref{align:ch07-7.115}，对所有 $\tau\in(0,T)\setminus N$ 与 $j\in\mathbb{N}$，得
\begin{align}
\begin{aligned}
\int_{\mathbb{R}}U_j^*(x,\tau)\zeta(x,\tau)\,\mathrm{d}x+\langle\sigma(\cdot,\tau),\zeta(\cdot,\tau)\rangle
\leqslant{}&\iint_{S_\tau}\{U_j^*\partial_t\zeta+\Lambda_{U_j}^*\partial_x\zeta\}\,\mathrm{d}x\,\mathrm{d}t\\
&+\int_0^\tau\langle\sigma(\cdot,t),\partial_t\zeta(\cdot,t)\rangle\,\mathrm{d}t
+\int_{\mathbb{R}}U_j(u_{0r})(x)\zeta(x,0)\,\mathrm{d}x+\langle u_{0s},\zeta(\cdot,0)\rangle.
\end{aligned} \label{align:ch07-7.116}
\end{align}
因对所有 $y\geqslant0$，$0\leqslant U_j(y)\leqslant y\chi_{[j,\infty)}(y)$ 且 $|\Lambda_{U_j}(y)|\leqslant My\chi_{[j,\infty)}(y)$（见\eqref{align:ch07-7.1}），故 $|U_j^*|\leqslant u_b$ 且 $|\Lambda_{U_j}^*|\leqslant Mu_b$（$j\in\mathbb{N}$）。因 $u_b\in L^\infty(0,T;L^1(\mathbb{R}))$，显然
\begin{align}
\lim_{j\to\infty}U_j^*=\lim_{j\to\infty}\Lambda_{U_j}^*=0\quad\text{在 }S\text{ 中几乎处处且在 }L^1(S), \label{align:ch07-7.117}
\end{align}
从而在\eqref{align:ch07-7.116}中令 $j\to\infty$，对所有 $\tau\in(0,T)\setminus N$ 与 $\zeta\in C^1([0,T];C^1_c(\mathbb{R}))$ 得
\begin{align}
\langle\sigma(\cdot,\tau),\zeta(\cdot,\tau)\rangle
\leqslant\int_0^\tau\langle\sigma(\cdot,t),\partial_t\zeta(\cdot,t)\rangle\,\mathrm{d}t+\langle u_{0s},\zeta(\cdot,0)\rangle. \label{align:ch07-7.118}
\end{align}
在\eqref{align:ch07-7.118}中取 $\zeta(x,t)=\rho(x)$，得对任意非负 $\rho\in C^1_c(\mathbb{R})$，$\langle\sigma(\cdot,\tau),\rho\rangle\leqslant\langle u_{0s},\rho\rangle$，故 $\sigma(\cdot,\tau)\ll u_{0s}$（回忆 $\sigma$ 非负）。于是对任意 $\tau\in(0,T)\setminus N$，$\sigma(\cdot,\tau)$ 关于 Lebesgue 测度奇异。另一方面，因 $u_b\in L^\infty(0,T;L^1(\mathbb{R}))$，测度 $E\mapsto\iint_E u_b\,\mathrm{d}x\,\mathrm{d}t$（$E\in\mathcal{B}(\mathbb{R}^2)\cap S$）关于 Lebesgue 测度绝对连续。由 Lebesgue 分解的唯一性（见\cref{theorem:ch01-Lebesgue分解}）得结论。
\end{proof}

\begin{proposition}\label{proposition:ch07-divcurl}\label{proposition:ch07-7.5.6}
设 (A0)--(A1) 成立。设 $\nu\in\mathcal{Y}^+(S;\mathbb{R})$ 由\cref{proposition:ch07-7.5.2}给出，$u\in L^\infty_{w^*}(0,T;\mathscr{R}_f(\mathbb{R}))$ 由\eqref{align:ch07-7.108}定义。则
\begin{align}
\phi^*(x,t)=\phi(u_r(x,t))\quad\text{对几乎所有 }(x,t)\in S \label{align:ch07-7.119}
\end{align}
其中 $\phi^*$ 由\eqref{align:ch07-7.4}定义。
\end{proposition}

\begin{proof}
设 $U,V\in C^2([0,\infty))\cap L^\infty(0,\infty)$ 满足\eqref{align:ch07-7.58}，并设 $\Lambda_U,\Lambda_V\in L^\infty(0,\infty)$（见\eqref{align:ch07-7.83}）。由\cref{lemma:ch07-7.4.8}，对任意 $n\in\mathbb{N}$，存在 $\sigma_n,\tau_n\in\mathscr{R}_f(S)$，使对所有 $\zeta\in C^1_c(S)$，
\begin{align}
\iint_S\{U(u_n)\partial_t\zeta+\Lambda_U(u_n)\partial_x\zeta\}\,\mathrm{d}x\,\mathrm{d}t=\langle\sigma_n,\zeta\rangle, \label{align:ch07-7.120a}
\end{align}
\begin{align}
\iint_S\{V(u_n)\partial_t\zeta+\Lambda_V(u_n)\partial_x\zeta\}\,\mathrm{d}x\,\mathrm{d}t=\langle\tau_n,\zeta\rangle, \label{align:ch07-7.120b}
\end{align}
其中 $u_n$ 为\cref{proposition:ch07-7.4.7}给出的 (Pn) 的熵解。

设 $A\Subset S$ 为有界开集，定义 $Y_n,Z_n:A\to\mathbb{R}^2$ 为
\begin{align*}
Y_n:=(\Lambda_U(u_n),U(u_n)),\qquad Z_n:=(V(u_n),-\Lambda_V(u_n)).
\end{align*}
由\eqref{align:ch07-7.120a}得
\begin{align}
\operatorname{div}Y_n=-\sigma_n,\qquad \operatorname{curl}Z_n=-\tau_n\quad\text{在 }\mathcal{D}^*(A). \label{align:ch07-7.121}
\end{align}
因 $U,\Lambda_U,V,\Lambda_V$ 在 $(0,\infty)$ 上有界，序列 $\{U(u_n)\}$、$\{\Lambda_U(u_n)\}$、$\{V(u_n)\}$、$\{\Lambda_V(u_n)\}$ 在 $L^\infty(A)$ 中有界，从而在 $L^1(A)$ 中有界且一致可积。由\cref{proposition:ch05-5.4.10}，在 $L^1(A)$ 中
\begin{align*}
U(u_n)\rightharpoonup U^*:=\int_{[0,\infty)}U(y)\,\mathrm{d}\nu(y),\quad
\Lambda_U(u_n)\rightharpoonup\Lambda_U^*:=\int_{[0,\infty)}\Lambda_U(y)\,\mathrm{d}\nu(y),
\end{align*}
\begin{align*}
V(u_n)\rightharpoonup V^*:=\int_{[0,\infty)}V(y)\,\mathrm{d}\nu(y),\quad
\Lambda_V(u_n)\rightharpoonup\Lambda_V^*:=\int_{[0,\infty)}\Lambda_V(y)\,\mathrm{d}\nu(y).
\end{align*}
同理，$\{U(u_n)\}$、$\{\Lambda_U(u_n)\}$、$\{V(u_n)\}$、$\{\Lambda_V(u_n)\}$ 有界，从而在 $L^2(A)$ 中弱收敛。于是当 $n\to\infty$ 时，
\begin{align}
Y_n\rightharpoonup Y^*:=(\Lambda_U^*,U^*),\quad
Z_n\rightharpoonup Z^*:=(V^*,-\Lambda_V^*)\quad\text{在 }L^2(A;\mathbb{R}^2). \label{align:ch07-7.122}
\end{align}
类似地，
\begin{align}
Y_n\cdot Z_n:=\Lambda_U(u_n)V(u_n)-\Lambda_V(u_n)U(u_n)
\rightharpoonup\int_{[0,\infty)}[\Lambda_U(y)V(y)-\Lambda_V(y)U(y)]\,\mathrm{d}\nu(y)
\quad\text{在 }L^2(A). \label{align:ch07-7.123}
\end{align}
由\eqref{align:ch07-7.85}（应用于 $\sigma_n$ 与 $\tau_n$）与\eqref{align:ch07-7.121}，$\{\operatorname{div}Y_n\}$ 与 $\{\operatorname{curl}Z_n\}$ 在 $H^{-1}(A)$ 中预紧。于是由\eqref{align:ch07-7.122}与 div--curl 引理（Murat--Tartar）得
\begin{align}
Y_n\cdot Z_n\to Y^*\cdot Z^*=\Lambda_U^*V^*-\Lambda_V^*U^*\quad\text{在 }\mathcal{D}^*(A) \label{align:ch07-7.124}
\end{align}
（例如见文献~\cite{ref46} 推论 1.3.1 与定理 5.2.1）。由\eqref{align:ch07-7.123}--\eqref{align:ch07-7.124}得
\begin{align}
\int_{[0,\infty)}[\Lambda_U(y)-\Lambda_U^*]V(y)\,\mathrm{d}\nu(y)
=\int_{[0,\infty)}[U(y)-U^*]\Lambda_V(y)\,\mathrm{d}\nu(y)
\quad\text{在 }A\text{ 中几乎处处}. \label{align:ch07-7.125}
\end{align}

设 $U$ 如上，$U'>0$ 在 $(0,\infty)$。由标准逼近论证可取 $V(y)=|U(y)-U^*|$，适当选取 $\vartheta_V$ 得
\begin{align*}
\Lambda_V(y)=\int_0^y V'(z)\phi'(z)\,\mathrm{d}z+\vartheta_V
=\operatorname{sgn}(U(y)-U^*)[\Lambda_U(y)-\Lambda_U(U^{-1}(U^*))],
\end{align*}
故由\eqref{align:ch07-7.125}，
\begin{align}
[\Lambda_U^*-\Lambda_U(U^{-1}(U^*))]\int_{[0,\infty)}|U(y)-U^*|\,\mathrm{d}\nu(y)=0
\quad\text{在 }A\text{ 中几乎处处}. \label{align:ch07-7.126}
\end{align}
设 $U_p\in C^2([0,\infty))\cap L^\infty(0,\infty)$，$U_p(0)=0$（$p\in\mathbb{N}$），满足\eqref{align:ch07-7.58}及
\begin{align}
0<U_p'\leqslant U_{p+1}'\leqslant1,\quad \lim_{p\to\infty}U_p'(y)=1\ \text{在 }[0,\infty), \label{align:ch07-7.127}
\end{align}
从而 $\lim_{p\to\infty}U_p(y)=y$（对所有 $y\in[0,\infty)$）。由单调收敛定理、\eqref{align:ch07-7.95}与\eqref{align:ch07-7.109a}得
\begin{align}
\lim_{p\to\infty}U_p^*=\lim_{p\to\infty}\int_{[0,\infty)}U_p(y)\,\mathrm{d}\nu(y)=u_r\quad\text{在 }A\text{ 中几乎处处}. \label{align:ch07-7.128}
\end{align}
我们证明
\begin{align}
\lim_{p\to\infty}[\Lambda_{U_p}^*-\Lambda_{U_p}(U_p^{-1}(U_p^*))]=\phi^*-\phi(u_r)\quad\text{在 }A\text{ 中几乎处处}. \label{align:ch07-7.129}
\end{align}
利用\eqref{align:ch07-7.128}--\eqref{align:ch07-7.129}即可完成证明。事实上，由\eqref{align:ch07-7.128}与控制收敛定理，对几乎所有 $(x,t)\in A$，
\begin{align}
\lim_{p\to\infty}\int_{[0,\infty)}|U_p(y)-U_p^*(x,t)|\,\mathrm{d}\nu_{(x,t)}(y)
=\int_{[0,\infty)}|y-u_r(x,t)|\,\mathrm{d}\nu_{(x,t)}(y). \label{align:ch07-7.130}
\end{align}
在\eqref{align:ch07-7.126}（取 $U=U_p$）中令 $p\to\infty$，由\eqref{align:ch07-7.129}--\eqref{align:ch07-7.130}，对几乎所有 $(x,t)\in A$ 得
\begin{align*}
[\phi^*(x,t)-\phi(u_r)(x,t)]\int_{[0,\infty)}|y-u_r(x,t)|\,\mathrm{d}\nu_{(x,t)}(y)=0.
\end{align*}
由 $A$ 任意，上式蕴含\eqref{align:ch07-7.119}。

余下需证\eqref{align:ch07-7.129}。注意
\begin{align}
\Lambda_{U_p}^*-\Lambda_{U_p}(U_p^{-1}(U_p^*))
=\int_{[0,\infty)}\left(\int_0^y U_p'(z)\phi'(z)\,\mathrm{d}z\right)\mathrm{d}\nu(y)
-\int_0^{U_p^{-1}(U_p^*)}U_p'(z)\phi'(z)\,\mathrm{d}z. \label{align:ch07-7.131}
\end{align}
由\eqref{align:ch07-7.1}与\eqref{align:ch07-7.127}，$U_p'(y)\to1$ 且 $|U_p'(y)\phi'(y)|\leqslant M$（对所有 $y\geqslant0$），由控制收敛定理，
\begin{align}
\lim_{p\to\infty}\Lambda_{U_p}^*(x,t)=\phi^*(x,t)\quad\text{对几乎所有 }(x,t)\in A. \label{align:ch07-7.132}
\end{align}
另一方面，
\begin{align}
\int_0^{U_p^{-1}(U_p^*(x,t))}U_p'(z)\phi'(z)\,\mathrm{d}z-\phi(u_r)(x,t)
=\int_0^{u_r(x,t)}[U_p'(z)-1]\phi'(z)\,\mathrm{d}z
+\int_{u_r(x,t)}^{U_p^{-1}(U_p^*(x,t))}U_p'(z)\phi'(z)\,\mathrm{d}z. \label{align:ch07-7.133}
\end{align}
如\eqref{align:ch07-7.132}那样论证，
\begin{align}
\lim_{p\to\infty}\int_0^{u_r(x,t)}[U_p'(z)-1]\phi'(z)\,\mathrm{d}z=0
\quad\text{对几乎所有 }(x,t)\in A. \label{align:ch07-7.134}
\end{align}
关于\eqref{align:ch07-7.133}右端第二项，对某 $\delta>0$，
\begin{align*}
\left|\int_{u_r(x,t)}^{U_p^{-1}(U_p^*(x,t))}U_p'(z)\phi'(z)\,\mathrm{d}z\right|
\leqslant M|u_r(x,t)-U_p^{-1}(U_p^*(x,t))|
\end{align*}
对充分大的 $p$（见\eqref{align:ch07-7.127}）。由\eqref{align:ch07-7.128}令 $p\to\infty$，得
\begin{align}
\lim_{p\to\infty}\int_{u_r(x,t)}^{U_p^{-1}(U_p^*(x,t))}U_p'(z)\phi'(z)\,\mathrm{d}z=0
\quad\text{对几乎所有 }(x,t)\in A. \label{align:ch07-7.135}
\end{align}
由\eqref{align:ch07-7.133}--\eqref{align:ch07-7.135}得
\begin{align}
\int_0^{U_p^{-1}(U_p^*(x,t))}U_p'(z)\phi'(z)\,\mathrm{d}z\to\phi(u_r)(x,t)
\quad\text{对几乎所有 }(x,t)\in A. \label{align:ch07-7.136}
\end{align}
由\eqref{align:ch07-7.131}、\eqref{align:ch07-7.132}与\eqref{align:ch07-7.136}得\eqref{align:ch07-7.129}。证毕。
\end{proof}

\begin{proof}[\cref{theorem:ch07-7.5.1}的证明]
设 $u\in L^\infty_{w^*}(0,T;\mathscr{R}_f(\Omega))$ 由\eqref{align:ch07-7.108}定义。则由\cref{proposition:ch07-7.5.2}--\cref{proposition:ch07-divcurl}得结论。
\end{proof}

现可证明\cref{theorem:ch07-7.3.1}。

\begin{proof}[\cref{theorem:ch07-7.3.1}的证明]
证明由\eqref{align:ch07-7.108}定义的 $u\in L^\infty_{w^*}(0,T;\mathscr{R}_f(\Omega))$ 满足等式
\begin{align}
\iint_S[u_r\partial_t\zeta+\phi(u_r)\partial_x\zeta]\,\mathrm{d}x\,\mathrm{d}t
+\int_0^T\langle u_s(\cdot,t),\partial_t\zeta(\cdot,t)\rangle\,\mathrm{d}t
=-\langle u_0,\zeta(\cdot,0)\rangle \label{align:ch07-7.137}
\end{align}
对所有 $\zeta\in C^1([0,T];C^1_c(\mathbb{R}))$、$\zeta(\cdot,T)=0$ 在 $\mathbb{R}$ 上成立，从而 $u$ 是 (P) 的解（见\cref{definition:ch07-熵解}）。设 $\{u_{n_k}\}\subseteq\{u_n\}$ 为\cref{theorem:ch07-7.5.1}给出的子列。对 $U(y)=y$ 与 $U(y)=\phi(y)$ 应用\cref{theorem:ch07-7.5.1}(iii)，得
\begin{align}
\lim_{k\to\infty}\iint_S u_{n_k}\partial_t\zeta\,\mathrm{d}x\,\mathrm{d}t
=\iint_S u_r\partial_t\zeta\,\mathrm{d}x\,\mathrm{d}t+\int_0^T\langle u_s(\cdot,t),\partial_t\zeta(\cdot,t)\rangle\,\mathrm{d}t, \label{align:ch07-7.138a}
\end{align}
\begin{align}
\lim_{k\to\infty}\iint_S\phi(u_{n_k})\partial_x\zeta\,\mathrm{d}x\,\mathrm{d}t
=\iint_S\phi(u_r)\partial_x\zeta\,\mathrm{d}x\,\mathrm{d}t \label{align:ch07-7.138b}
\end{align}
（见\eqref{align:ch07-7.89}、\eqref{align:ch07-7.4}与\eqref{align:ch07-7.119}）。由\eqref{align:ch07-7.45}中第一个收敛与\eqref{align:ch07-7.138a}，在\eqref{align:ch07-7.78}（取 $n=n_k$）中令 $k\to\infty$ 得\eqref{align:ch07-7.137}。不等式\eqref{align:ch07-7.5}类似地在\eqref{align:ch07-7.77}（取 $n=n_k$）中令 $k\to\infty$ 证明，对 $U(y)=E_l(y)$ 与 $U(y)=F_l(y)$ 应用\cref{theorem:ch07-7.5.1}(iii)，并如\cref{proposition:ch07-7.5.3}证明中那样论证 $E_l(u_{0n_k})\stackrel{*}{\rightharpoonup}E_l(u_{0r})+u_{0s}$ 在 $\mathscr{R}_f(\mathbb{R})$ 中。结论得证。
\end{proof}

为证\cref{theorem:ch07-7.3.2}，需要下列结果，它刻画 Young 测度 $\nu$ 的离散分解。

\begin{proposition}\label{proposition:ch07-7.5.7}
设 (A0)--(A2) 成立。设 $u\in L^\infty_{w^*}(0,T;\mathscr{R}_f(\mathbb{R}))$ 与 $\nu\in\mathcal{Y}^+(S;\mathbb{R})$ 由\cref{theorem:ch07-7.5.1}给出。则对几乎所有 $(x,t)\in S$，
\begin{align}
\nu_{(x,t)}=\delta_{u_r(x,t)}. \label{align:ch07-7.139}
\end{align}
\end{proposition}

\begin{proof}
固定 $(x,t)\in S$。若 $u_r(x,t)=0$，由\eqref{align:ch07-7.89}立即得\eqref{align:ch07-7.139}。设 $y_1\equiv u_r(x,t)>0$。对所有 $y_2>y_1$ 与充分大的 $p\in\mathbb{N}$，置
\begin{align*}
V_p(y):=p(y-y_1)\chi_{(y_1,y_1+1/p)}(y)+\chi_{[y_1+1/p,y_2)}(y)+p(y_2+1/p-y)\chi_{[y_2,y_2+1/p)}(y),
\end{align*}
\begin{align*}
\Lambda_{V_p}(y):=\int_0^y V_p'(z)\phi'(z)\,\mathrm{d}z\quad(y\in[0,\infty)).
\end{align*}
则对任意 $y\in[0,\infty)$，
\begin{align}
\lim_{p\to\infty}V_p(y)=\chi_{(y_1,y_2]}(y), \label{align:ch07-7.140a}
\end{align}
\begin{align}
\lim_{p\to\infty}\Lambda_{V_p}(y)=\phi'(y_1)\chi_{(y_1,y_2]}(y)+[\phi'(y_1)-\phi'(y_2)]\chi_{(y_2,\infty)}(y). \label{align:ch07-7.140b}
\end{align}
设 $\{U_p\}$ 为\cref{proposition:ch07-7.5.6}证明中所用序列。由标准逼近论证，可在等式\eqref{align:ch07-7.125}中取 $U=U_p$、$V=V_p$，得
\begin{align*}
\int_{[0,\infty)}[\Lambda_{U_p}(y)-\Lambda_{U_p}^*(x,t)]V_p(y)\,\mathrm{d}\nu_{(x,t)}(y)
=\int_{[0,\infty)}[U_p(y)-U_p^*(x,t)]\Lambda_{V_p}(y)\,\mathrm{d}\nu_{(x,t)}(y).
\end{align*}
由\eqref{align:ch07-7.128}、\eqref{align:ch07-7.132}与\eqref{align:ch07-7.119}，对所有 $y\geqslant0$ 有
\begin{align*}
\lim_{p\to\infty}U_p(y)-U_p^*(x,t)=y-y_1,
\end{align*}
\begin{align*}
\lim_{p\to\infty}\Lambda_{U_p}^*(x,t)-\Lambda_{U_p}(y)=\phi(y_1)-\phi(y).
\end{align*}
由\eqref{align:ch07-7.140a}--\eqref{align:ch07-7.140b}，在 $S$ 中几乎处处得
\begin{align*}
\lim_{p\to\infty}\int_{[0,\infty)}[\Lambda_{U_p}(y)-\Lambda_{U_p}^*]V_p(y)\,\mathrm{d}\nu(y)
=\int_{(y_1,y_2]}[\phi(y)-\phi(y_1)]\,\mathrm{d}\nu(y),
\end{align*}
\begin{align*}
\begin{aligned}
\lim_{p\to\infty}\int_{[0,\infty)}&[U_p(y)-U_p^*]\Lambda_{V_p}(y)\,\mathrm{d}\nu(y)\\
&=\int_{(y_1,y_2]}\phi'(y_1)(y-y_1)\,\mathrm{d}\nu(y)
+[\phi'(y_1)-\phi'(y_2)]\int_{(y_2,\infty)}(y-y_1)\,\mathrm{d}\nu(y),
\end{aligned}
\end{align*}
从而
\begin{align}
\int_{(y_1,y_2]}[\phi(y)-\phi(y_1)-\phi'(y_1)(y-y_1)]\,\mathrm{d}\nu(y)
=[\phi'(y_1)-\phi'(y_2)]\int_{(y_2,\infty)}(y-y_1)\,\mathrm{d}\nu(y). \label{align:ch07-7.142a}
\end{align}
类似地，对任意 $y_0\in(0,y_1)$，
\begin{align}
\begin{aligned}
\int_{(y_0,y_1)}&[\phi(y)-\phi(y_1)-\phi'(y_1)(y-y_1)]\,\mathrm{d}\nu_{(x,t)}(y)\\
&=[\phi'(y_1)-\phi'(y_0)]\int_{[0,y_0]}(y-y_1)\,\mathrm{d}\nu_{(x,t)}(y).
\end{aligned} \label{align:ch07-7.142b}
\end{align}
由假设 (A2)，$\phi'$ 至少在区间 $[y_0,y_1]$ 与 $[y_1,y_2]$ 之一上严格单调。设在 $[y_1,y_2]$ 上，另一情形类似。则由\eqref{align:ch07-7.142a}，对几乎所有 $(x,t)\in S$，
\begin{align*}
\int_{(y_1,y_2]}|\phi(y)-\phi(y_1)-\phi'(y_1)(y-y_1)|\,\mathrm{d}\nu_{(x,t)}(y)
+|\phi'(y_1)-\phi'(y_2)|\int_{(y_2,\infty)}|y_1-y|\,\mathrm{d}\nu_{(x,t)}(y)=0,
\end{align*}
故 $\operatorname{supp}\nu_{(x,t)}\subseteq[0,y_1]$。因 $\nu_{(x,t)}$ 为概率测度，
\begin{align*}
u_r(x,t)=\int_{[0,u_r(x,t)]}y\,\mathrm{d}\nu_{(x,t)}(y)
=\int_{[0,u_r(x,t)]}[y-u_r(x,t)]\,\mathrm{d}\nu_{(x,t)}(y)+u_r(x,t),
\end{align*}
从而
\begin{align*}
\int_{[0,u_r(x,t)]}|y-u_r(x,t)|\,\mathrm{d}\nu_{(x,t)}(y)=0.
\end{align*}
由上式得 $\operatorname{supp}\nu_{(x,t)}=\{u_r(x,t)\}$，因 $\nu_{(x,t)}$ 为概率测度，\eqref{align:ch07-7.139}得证。
\end{proof}

\begin{remark}\label{remark:ch07-7.5.2}
由\eqref{align:ch07-7.139}与\cref{proposition:ch05-5.4.1}，可能抽取子列（不重新标号），有
\begin{align}
u_{n_k}\to u_r\quad\text{在 }S\text{ 中几乎处处}. \label{align:ch07-7.143}
\end{align}
特别地，若 $\phi$ 有界，由控制收敛定理，对任意 $L>0$，$\phi(u_{n_k})\to\phi(u_r)$ 在 $L^1(R_L)$ 中。
\end{remark}

\begin{remark}\label{remark:ch07-7.5.3}
设 (A0)--(A1) 成立，且对某 $a>0$、$b>0$ 与 $(x,t)\in S$，$\phi'$ 在区间 $I_{a,b}=[u_r(x,t)-a,u_r(x,t)+b]$ 上为常数。此情形可证 $\operatorname{supp}\nu_{(x,t)}\subseteq I_{(x,t)}$，其中 $I_{(x,t)}$ 为 $\phi'=\phi'(u_r)$ 的最大区间（见文献~\cite{ref18} 命题 5.9）。
\end{remark}

\begin{proof}[\cref{theorem:ch07-7.3.2}的证明]
由\cref{theorem:ch07-7.3.1}，$u$ 是 (P) 的 Young 测度熵解，故不等式\eqref{align:ch07-7.5}成立。由\eqref{align:ch07-7.5}与\eqref{align:ch07-7.139}得\eqref{align:ch07-7.8}，结论得证。
\end{proof}

最后证明\cref{lemma:ch07-7.5.4}。

\begin{proof}[\cref{lemma:ch07-7.5.4}的证明]
由标准逼近论证，只需证等式\eqref{align:ch07-7.106a}。

(i) 先在 $U\in C^2([0,\infty))$、$U'\in L^\infty(0,\infty)$ 满足\eqref{align:ch07-7.58}与\eqref{align:ch07-7.91}的假设下证\eqref{align:ch07-7.106a}。设 $\rho\in C^2_c(\mathbb{R})$，$h\in C_c(0,T)$，$\operatorname{supp}\rho\subseteq I_L$（某 $L>0$）。设 $U_{n_k,\rho}$ 由\eqref{align:ch07-7.60}定义。由\eqref{align:ch07-7.104}得
\begin{align}
\lim_{k\to\infty}\int_0^T U_{n_k,\rho}(t)h(t)\,\mathrm{d}t
=\int_0^T U_\rho^*(t)h(t)\,\mathrm{d}t+M_U\int_0^T\langle\sigma(\cdot,t),\rho\rangle h(t)\,\mathrm{d}t. \label{align:ch07-7.144}
\end{align}
另一方面，由\cref{lemma:ch07-7.4.6}，序列 $\{U_{n_k,\rho}\}$ 在 $BV(0,T)$ 中有界，故存在子列（不重新标号）在 $L^1(0,T)$ 中收敛于 $U_\rho^*+M_U\langle\sigma(\cdot,\cdot),\rho\rangle$（见\eqref{align:ch07-7.144}）。因 $\{U(u_{n_k})\}$ 在 $L^\infty(0,T;L^1(I_L))$ 中有界且 $U^*\in L^\infty(0,T;L^1(I_L))$（见\cref{proposition:ch07-7.5.3}与不等式\eqref{align:ch07-7.101}），条件 $\rho\in C^2_c(\mathbb{R})$ 可放宽为 $\rho\in C_c(\mathbb{R})$。故此情形下\eqref{align:ch07-7.106a}成立。

(ii) 现对所有 $U\in C_b([0,\infty))$ 证\eqref{align:ch07-7.106a}（注意此情形 $M_U=0$）。置
\begin{align*}
U_p:[0,\infty)\to\mathbb{R},\quad U_p(y):=((U\chi_{[0,p]})*\theta_p)(y)\quad(y\geqslant0),
\end{align*}
其中 $\theta_p(y):=p\theta(py)$，$\theta\in C^\infty_c(\mathbb{R})$ 为标准磨光核（$p\in\mathbb{N}$）。则 $\{U_p\}\subseteq C^2_c([0,\infty))$，$U_p\to U$ 在 $[0,\infty)$ 的紧子集上一致收敛，且 $\|U_p\|_{L^\infty(\mathbb{R})}\leqslant\|U\|_{L^\infty(\mathbb{R})}$。此外 $U_p'\in L^\infty(0,\infty)$，$U_p$ 满足\eqref{align:ch07-7.58}与\eqref{align:ch07-7.91}（$M_{U_p}=0$），故由 (i) 得
\begin{align}
\lim_{k\to\infty}\int_0^T\left|\int_{\mathbb{R}}U_p(u_{n_k}(x,t))\rho(x)\,\mathrm{d}x-U_{p,\rho}^*(t)\right|\mathrm{d}t=0, \label{align:ch07-7.145}
\end{align}
其中 $U_p^*:=\int_{[0,\infty)}U_p(y)\,\mathrm{d}\nu(y)$，$U_{p,\rho}^*:=\int_{\mathbb{R}}U_p^*(x,t)\rho(x)\,\mathrm{d}x$（对所有 $\rho\in C_c(\mathbb{R})$，$p\in\mathbb{N}$）。现注意
\begin{align}
\begin{aligned}
\int_0^T&\left|\int_{\mathbb{R}}U(u_{n_k}(x,t))\rho(x)\,\mathrm{d}x-U_\rho^*(t)\right|\mathrm{d}t\\
&\leqslant\iint_S|U(u_{n_k}(x,t))-U_p(u_{n_k}(x,t))|\,|\rho(x)|\,\mathrm{d}x\,\mathrm{d}t\\
&\quad+\int_0^T\left|\int_{\mathbb{R}}U_p(u_{n_k}(x,t))\rho(x)\,\mathrm{d}x-U_{p,\rho}^*(t)\right|\mathrm{d}t
+\int_0^T|U_\rho^*(t)-U_{p,\rho}^*(t)|\,\mathrm{d}t.
\end{aligned} \label{align:ch07-7.146}
\end{align}
关于\eqref{align:ch07-7.146}右端第一项，对所有 $p\in\mathbb{N}$ 与 $M>0$，由 Chebyshev 不等式与\eqref{align:ch07-7.71}中第二个不等式，
\begin{align}
\begin{aligned}
\iint_S&|U(u_{n_k})-U_p(u_{n_k})|\,|\rho|\,\mathrm{d}x\,\mathrm{d}t\\
&\leqslant\iint_{\{u_{n_k}\leqslant M\}}|U(u_{n_k})-U_p(u_{n_k})|\,|\rho|\,\mathrm{d}x\,\mathrm{d}t
+\iint_{\{u_{n_k}>M\}}|U(u_{n_k})-U_p(u_{n_k})|\,|\rho|\,\mathrm{d}x\,\mathrm{d}t\\
&\leqslant\|\rho\|_{L^\infty(\mathbb{R})}T\left\{\lambda(\operatorname{supp}\rho)\|U-U_p\|_{L^\infty(0,M)}+\frac{2}{M}\|u_0\|_{\mathscr{R}_f(\mathbb{R})}\|U\|_{L^\infty(\mathbb{R})}\right\}.
\end{aligned} \label{align:ch07-7.147}
\end{align}
关于\eqref{align:ch07-7.146}右端第三项，
\begin{align}
\begin{aligned}
\int_0^T|U_\rho^*(t)-U_{p,\rho}^*(t)|\,\mathrm{d}t
&\leqslant\iint_S|U^*-U_p^*|\,|\rho|\,\mathrm{d}x\,\mathrm{d}t\\
&\leqslant\|\rho\|_{L^\infty(\mathbb{R})}\left\{\lambda(\operatorname{supp}\rho)T\|U-U_p\|_{L^\infty(0,M)}
+2\|U\|_{L^\infty(\mathbb{R})}\iint_{\operatorname{supp}\rho\times(0,T)}\nu_{(x,t)}(\{y>M\})\,\mathrm{d}x\,\mathrm{d}t\right\}.
\end{aligned} \label{align:ch07-7.148}
\end{align}
由\eqref{align:ch07-7.145}--\eqref{align:ch07-7.148}，对所有 $\rho\in C_c(\mathbb{R})$、$p\in\mathbb{N}$ 与 $M>0$，
\begin{align*}
\begin{aligned}
\limsup_{k\to\infty}\int_0^T\mathrm{d}t\,&\left|\int_{\mathbb{R}}U(u_{n_k}(x,t))\rho(x)\,\mathrm{d}x-U_\rho^*(t)\right|\\
&\leqslant2\|\rho\|_{L^\infty(\mathbb{R})}\Big\{\lambda(\operatorname{supp}\rho)T\|U-U_p\|_{L^\infty(0,M)}
+\|U\|_{L^\infty(\mathbb{R})}\left[\frac{T}{M}+\iint_{\operatorname{supp}\rho\times(0,T)}\nu_{(x,t)}(\{y>M\})\,\mathrm{d}x\,\mathrm{d}t\right]\Big\}.
\end{aligned}
\end{align*}
因 $U_p\to U$ 在 $[0,\infty)$ 紧集上一致收敛，令 $p\to\infty$ 得
\begin{align}
\limsup_{k\to\infty}\int_0^T\left|\int_{\mathbb{R}}U(u_{n_k})\rho(x)\,\mathrm{d}x-U_\rho^*(t)\right|\mathrm{d}t
\leqslant2\|\rho\|_{L^\infty(\mathbb{R})}\|U\|_{L^\infty(\mathbb{R})}\left\{\frac{T}{M}+\iint_{\operatorname{supp}\rho\times(0,T)}\nu_{(x,t)}(\{y>M\})\,\mathrm{d}x\,\mathrm{d}t\right\}. \label{align:ch07-7.149}
\end{align}
因 $\nu_{(x,t)}$ 为概率测度，由控制收敛定理显然
\begin{align*}
\lim_{M\to\infty}\iint_{\operatorname{supp}\rho\times(0,T)}\nu_{(x,t)}(\{y>M\})\,\mathrm{d}x\,\mathrm{d}t=0.
\end{align*}
在\eqref{align:ch07-7.149}中令 $M\to\infty$，即得\eqref{align:ch07-7.106a}。

(iii) 现设 $U\in C([0,\infty))$ 为任意满足\eqref{align:ch07-7.91}的函数。如\cref{proposition:ch07-7.5.3}证明中那样，设 $f_{j1},f_{j2}\in C^2([0,\infty))$ 满足\eqref{align:ch07-7.100}（$j\in\mathbb{N}$，$j\geqslant[y_\epsilon]+1$）。则
\begin{align*}
U=U_{j1}+U_{j2},\quad U^*=U_{j1}^*+U_{j2}^*,\quad U_\rho^*=U_{j1,\rho}^*+U_{j2,\rho}^*,
\end{align*}
其中 $U_{jl}:=f_{jl}U$，$U_{jl}^*:=\int_{[0,\infty)}U_{jl}(y)\,\mathrm{d}\nu(y)$，$U_{jl,\rho}^*:=\int_{\mathbb{R}}U_{jl}^*(x,t)\rho(x)\,\mathrm{d}x$（$\rho\in C_c(\mathbb{R})$，$l=1,2$）。因此
\begin{align}
\begin{aligned}
\int_0^T&\left|\int_{\mathbb{R}}U(u_{n_k}(x,t))\rho(x)\,\mathrm{d}x-U_\rho^*(t)-M_U\langle\sigma(\cdot,t),\rho\rangle\right|\mathrm{d}t\\
&\leqslant\int_0^T\left|\int_{\mathbb{R}}U_{j1}(u_{n_k}(x,t))\rho(x)\,\mathrm{d}x-U_{j1,\rho}^*(t)\right|\mathrm{d}t\\
&\quad+\int_0^T\left|\int_{\mathbb{R}}U_{j2}(u_{n_k}(x,t))\rho(x)\,\mathrm{d}x-U_{j2,\rho}^*(t)-M_U\langle\sigma(\cdot,t),\rho\rangle\right|\mathrm{d}t=:I_{jk1}+I_{jk2}.
\end{aligned} \label{align:ch07-7.151}
\end{align}
因 $U_{j1}\in C_b([0,\infty))$ 从而 $M_{U_{j1}}=0$，由 (ii) 与\eqref{align:ch07-7.106a}得
\begin{align}
\lim_{k\to\infty}I_{jk1}=0\quad(j\in\mathbb{N},\ j\geqslant[y_\epsilon]+1). \label{align:ch07-7.152}
\end{align}
另一方面，
\begin{align}
\begin{aligned}
I_{jk2}\leqslant{}&\int_0^T\left|\int_{\mathbb{R}}[U_{j2}(u_{n_k}(x,t))-M_U f_{j2}(u_{n_k}(x,t))u_{n_k}(x,t)]\rho(x)\,\mathrm{d}x\right|\mathrm{d}t\\
&+M_U\int_0^T\mathrm{d}t\left|\int_{\mathbb{R}}f_{j2}(u_{n_k}(x,t))u_{n_k}(x,t)\rho(x)\,\mathrm{d}x-\int_{\mathbb{R}}F_{j2}^*(x,t)\rho(x)\,\mathrm{d}x-\langle\sigma(\cdot,t),\rho\rangle\right|\\
&+\iint_S|U_{j2}^*-M_U F_{j2}^*|(x,t)\,|\rho(x)|\,\mathrm{d}x\,\mathrm{d}t,
\end{aligned} \label{align:ch07-7.153}
\end{align}
其中 $F_{j2}^*:=\int_{[0,\infty)}f_{j2}(y)y\,\mathrm{d}\nu(y)$。

分别处理\eqref{align:ch07-7.153}右端三项。由\eqref{align:ch07-7.99}与\eqref{align:ch07-7.71}中第二个不等式，
\begin{align}
\int_0^T\left|\int_{\mathbb{R}}[U_{j2}(u_{n_k}(x,t))-M_U f_{j2}(u_{n_k}(x,t))u_{n_k}(x,t)]\rho(x)\,\mathrm{d}x\right|\mathrm{d}t
\leqslant\epsilon\|\rho\|_{L^\infty(\mathbb{R})}T\|u_0\|_{\mathscr{R}_f(\mathbb{R})}. \label{align:ch07-7.154}
\end{align}
此外，映射 $y\mapsto f_{j2}(y)y$ 属于 $C^2([0,\infty))$，导数有界，满足\eqref{align:ch07-7.58}与\eqref{align:ch07-7.91}（$M_U=1$）。故由 (i)，
\begin{align}
\lim_{k\to\infty}\int_0^T\mathrm{d}t\left|\int_{\mathbb{R}}f_{j2}(u_{n_k}(x,t))u_{n_k}(x,t)\rho(x)\,\mathrm{d}x-\int_{\mathbb{R}}F_{j2}^*(x,t)\rho(x)\,\mathrm{d}x-\langle\sigma(\cdot,t),\rho\rangle\right|=0. \label{align:ch07-7.155}
\end{align}
现证
\begin{align}
\lim_{j\to\infty}\iint_S|U_{j2}^*-M_U F_{j2}^*|(x,t)\,|\rho(x)|\,\mathrm{d}x\,\mathrm{d}t=0. \label{align:ch07-7.156}
\end{align}
对几乎所有 $(x,t)\in S$，
\begin{align*}
|U_{j2}^*(x,t)-M_U F_{j2}^*(x,t)|
\leqslant\int_{[0,\infty)}|U_{j2}(y)-M_U f_{j2}(y)y|\,\mathrm{d}\nu_{(x,t)}(y)
=\int_{[0,\infty)}f_{j2}(y)|U(y)-M_U y|\,\mathrm{d}\nu_{(x,t)}(y)
\leqslant\epsilon\int_{[j,\infty)}y\,\mathrm{d}\nu_{(x,t)}(y)\leqslant\epsilon u_b(x,t)
\end{align*}
（见\eqref{align:ch07-7.95}）。因 $u_b\in L^\infty(0,T;L^1(\mathbb{R}))$ 且对几乎所有 $(x,t)\in S$ 有 $\lim_{j\to\infty}\int_{[j,\infty)}y\,\mathrm{d}\nu_{(x,t)}(y)=0$，由控制收敛定理得等式\eqref{align:ch07-7.156}。

由\eqref{align:ch07-7.153}--\eqref{align:ch07-7.156}，
\begin{align*}
\limsup_{j\to\infty}\left(\limsup_{k\to\infty}I_{jk2}\right)\leqslant\epsilon\|\rho\|_{L^\infty(\mathbb{R})}T\|u_0\|_{\mathscr{R}_f(\mathbb{R})},
\end{align*}
故由\eqref{align:ch07-7.151}--\eqref{align:ch07-7.152}，
\begin{align*}
\limsup_{k\to\infty}\int_0^T\left|\int_{\mathbb{R}}[U(u_{n_k})-U^*](x,t)\rho(x)\,\mathrm{d}x-M_U\langle\sigma(\cdot,t),\rho\rangle\right|\mathrm{d}t
\leqslant\epsilon T\|u_0\|_{\mathscr{R}_f(\mathbb{R})}\|\rho\|_{L^\infty(\mathbb{R})}.
\end{align*}
由 $\epsilon$ 任意，上式给出\eqref{align:ch07-7.106a}。证毕。
\end{proof}

\section{正则性结果的证明}\label{section:ch07-正则证明}

\begin{proof}[\cref{proposition:ch07-7.1.2}的证明]
(i) 只证\eqref{align:ch07-7.15}，\eqref{align:ch07-7.16}类似。设 $\widetilde\zeta\in C^1([0,T];C^1_c(\mathbb{R}))$，$\widetilde\zeta(\cdot,T)=0$ 且 $\widetilde\zeta\geqslant0$。在\eqref{align:ch07-7.9}中令 $l\to\infty$，易得
\begin{align}
\int_0^T\langle u_s(\cdot,t),\partial_t\widetilde\zeta(\cdot,t)\rangle\,\mathrm{d}t\geqslant-\langle u_{0s},\widetilde\zeta(\cdot,0)\rangle. \label{align:ch07-7.157}
\end{align}
由 $C([0,T];C_c(\mathbb{R}))$ 的可分性，存在零测集 $N\subseteq(0,T)$，使对所有 $\zeta\in C([0,T];C_c(\mathbb{R}))$ 与 $\tau\in(0,T)\setminus N$，有
\begin{align}
\lim_{h\to0^+}\frac1h\int_\tau^{\tau+h}\langle u_s(\cdot,t),\zeta(\cdot,t)\rangle\,\mathrm{d}t
=\langle u_s(\cdot,\tau),\zeta(\cdot,\tau)\rangle. \label{align:ch07-7.158}
\end{align}
设 $t_1,t_2\in(0,T)\setminus N$，$0<t_1<t_2<T$。由标准逼近论证，可在\eqref{align:ch07-7.157}中取 $\widetilde\zeta(x,t)=g_h(t)\zeta(x,t)$，其中
\begin{align*}
g_h(t):=\frac1h(t-t_1)\chi_{[t_1,t_1+h]}(t)+\chi_{(t_1+h,t_2)}(t)+\frac1h(t_2+h-t)\chi_{[t_2,t_2+h]}(t),
\end{align*}
$h\in(0,\min\{t_2-t_1,T-t_2\})$。令 $h\to0^+$ 并利用\eqref{align:ch07-7.158}，得
\begin{align}
\langle u_s(\cdot,t_2),\zeta(\cdot,t_2)\rangle
\leqslant\int_{t_1}^{t_2}\langle u_s(\cdot,t),\partial_t\zeta(\cdot,t)\rangle\,\mathrm{d}t
+\langle u_s(\cdot,t_1),\zeta(\cdot,t_1)\rangle. \label{align:ch07-7.159}
\end{align}
在\eqref{align:ch07-7.159}中取 $\zeta(x,t)=\rho(x)$（$\rho\in C_c(\mathbb{R})$，$\rho\geqslant0$），即得\eqref{align:ch07-7.15}。

(ii) 如 (i) 那样论证，存在零测集 $N\subseteq(0,T)$，使对所有 $\rho\in C_c(\mathbb{R})$ 与 $\tau\in(0,T)\setminus N$，有
\begin{align}
\lim_{h\to0^+}\frac1h\int_\tau^{\tau+h}\langle u(\cdot,t),\rho\rangle\,\mathrm{d}t=\langle u(\cdot,\tau),\rho\rangle. \label{align:ch07-7.160}
\end{align}
置 $\Omega_j:=[-j-1,-j]\cup[j,j+1]$，$Q_{j,\tau}:=\Omega_j\times(0,\tau)$（$\tau\in(0,T)$，$j\in\mathbb{N}$）。设 $\{\rho_j\}\subseteq C^1_c(\mathbb{R})$ 满足 $\rho_j=1$ 在 $[-j,j]$，$\operatorname{supp}\rho_j\subseteq[-j-1,j+1]$，$0\leqslant\rho_j\leqslant1$ 且 $|\rho_j'|\leqslant2$。在\eqref{align:ch07-7.7}中取 $\zeta(x,t)=f_h(t)\rho_j(x)$，其中
\begin{align*}
f_h(t):=\chi_{[0,\tau)}(t)+\frac1h(\tau+h-t)\chi_{[\tau,\tau+h]}(t),
\end{align*}
$h$ 充分小。则
\begin{align*}
\iint_{Q_{j,\tau}}f_h(t)\phi(u_r)(x,t)\rho_j'(x)\,\mathrm{d}x\,\mathrm{d}t
=\frac1h\int_\tau^{\tau+h}\langle u(\cdot,t),\rho_j\rangle\,\mathrm{d}t-\langle u_0,\rho_j\rangle.
\end{align*}
令 $h\to0^+$ 并利用\eqref{align:ch07-7.160}，得
\begin{align}
\iint_{Q_{j,\tau}}\phi(u_r)(x,t)\rho_j'(x)\,\mathrm{d}x\,\mathrm{d}t
=\langle u(\cdot,\tau),\rho_j\rangle-\langle u_0,\rho_j\rangle. \label{align:ch07-7.161}
\end{align}
因 $|\phi(u_r)|\leqslant Mu_r$（见\eqref{align:ch07-7.1}）且 $u_r\in L^\infty(0,T;L^1(\mathbb{R}))$，在\eqref{align:ch07-7.161}中令 $j\to\infty$，即得\eqref{align:ch07-7.17}。
\end{proof}

\begin{proof}[\cref{proposition:ch07-7.2.1}的证明]
我们证明如下结论。

\paragraph{结论 1.}对几乎所有 $t\in(0,T)$，映射 $x\mapsto Y(x,t):=\int_0^t\phi(u_r(x,s))\,\mathrm{d}s$ 属于 $BV(\mathbb{R})$；且对任意 $a,b\in\mathbb{R}$，$a\leqslant b$，有
\begin{align}
u(\cdot,t)([a,b])-u_0([a,b])=Y(a^-,0,t)-Y(b^+,0,t). \label{align:ch07-7.162}
\end{align}

取 $a=b=x_0$。由\eqref{align:ch07-7.162}，对几乎所有 $0\leqslant t\leqslant T$，
\begin{align*}
u_s(\cdot,t)(\{x_0\})=u_{0s}(\{x_0\})+Y(x_0^-,0,t)-Y(x_0^+,0,t)
\geqslant u_{0s}(\{x_0\})-2\|\phi\|_\infty t,
\end{align*}
故当 $t\in(0,u_{0s}(\{x_0\})/(2\|\phi\|_\infty))$ 时 $u_s(\cdot,t)(\{x_0\})>0$。由此得\eqref{align:ch07-7.19}。

余下证结论。如\eqref{align:ch07-7.161}那样论证，得
\begin{align}
\int_\mathbb{R}Y(x,\tau)\rho'(x)\,\mathrm{d}x=\langle u(\cdot,\tau),\rho\rangle-\langle u_0,\rho\rangle. \label{align:ch07-7.163}
\end{align}
故分布导数 $\partial_xY(x,t)$ 属于 $\mathscr{R}_f(\mathbb{R})$。另一方面，由\eqref{align:ch07-7.1}，$|Y(\cdot,t)|\leqslant M\int_0^t u_r(\cdot,s)\,\mathrm{d}s\in L^1(\mathbb{R})$，故对几乎所有 $t\in(0,T)$，$Y(\cdot,t)\in BV(\mathbb{R})$。

为证\eqref{align:ch07-7.162}，由标准正则化论证可在\eqref{align:ch07-7.163}中取 $\rho=\rho_j$，
\begin{align*}
\rho_j(x):=j\left(x-a+\frac1j\right)\chi_{[a-1/j,a]}(x)+\chi_{(a,b)}(x)+j\left(b+\frac1j-x\right)\chi_{[b,b+1/j]}(x),
\end{align*}
其中 $a,b\in\mathbb{R}$，$a<b$，$j\in\mathbb{N}$。于是
\begin{align}
\langle u(\cdot,\tau),\rho_j\rangle-\langle u_0,\rho_j\rangle
=j\int_{a-1/j}^a Y(x,\tau)\,\mathrm{d}x-j\int_b^{b+1/j}Y(x,\tau)\,\mathrm{d}x. \label{align:ch07-7.164}
\end{align}
在\eqref{align:ch07-7.164}中令 $j\to\infty$，即得\eqref{align:ch07-7.162}。结论得证。
\end{proof}

\begin{remark}\label{remark:ch07-7.6.1}
\cref{proposition:ch07-7.2.1}的证明事实上依赖于一维双曲守恒律的熵解与 Hamilton--Jacobi 方程的粘性解之间的对应（参见文献~\cite{ref21,ref20}）。
\end{remark}

不妨在 (A2') 成立且 $\phi''<0$、$\phi'>0$ 在 $(0,\infty)$、$\phi(0)=0$（从而 $\alpha\phi+\beta>0$，$\beta>0$，且 (A1) 亦满足；见\cref{remark:ch07-7.1.1}）的假设下证明\cref{theorem:ch07-7.3.3}与\cref{proposition:ch07-7.3.5}。考虑 (Pn) 的下列正则化问题（与\ref{section:ch07-逼近问题}中所用的 (P$n^\epsilon$) 不同）：
\begin{align*}
\begin{cases}
\partial_t y_n^\epsilon+\partial_x[\phi(y_n^\epsilon)]=\epsilon\,\partial_{xx}[\phi(y_n^\epsilon)] & \text{在 }S,\\
y_n^\epsilon=u_{0n}^\epsilon & \text{在 }\mathbb{R}\times\{0\},
\end{cases}
\qquad(Vn^\epsilon)
\end{align*}
其中 $\{u_{0n}^\epsilon\}$ 满足\eqref{align:ch07-7.46}--\eqref{align:ch07-7.47}。\ref{section:ch07-逼近问题}中为 (P$n^\epsilon$) 回顾的存在唯一与正则性结果、\cref{lemma:ch07-7.4.1}的先验估计及\cref{lemma:ch07-7.4.5}的收敛结果，对 (V$n^\epsilon$) 的解同样成立（见文献~\cite{ref63}）。族 $\{y_n^\epsilon\}$ 的若干性质是下列两个引理的内容。

\begin{lemma}\label{lemma:ch07-7.6.1}
设 (A2') 成立，$\phi(0)=0$，$\phi''<0$、$\phi'>0$ 在 $(0,\infty)$。设 $u_n$ 为\cref{proposition:ch07-7.4.7}给出的问题 (Pn) 的唯一熵解。则存在 (V$n^\epsilon$) 解的子列 $\{y_n^{\epsilon_j}\}$，使
\begin{align}
y_n^{\epsilon_j}\stackrel{*}{\rightharpoonup}u_n\quad\text{在 }L^\infty(S),\qquad
y_n^{\epsilon_j}\to u_n\quad\text{在 }L^1(R_L)\ \text{对所有 }L>0. \label{align:ch07-7.165}
\end{align}
\end{lemma}

\begin{proof}
由\eqref{align:ch07-7.66}与\eqref{align:ch07-7.69}，存在序列 $\{y_n^{\epsilon_j}\}$ 与 $y_n\in L^\infty(S)\cap L^\infty(0,T;L^1(\mathbb{R}))$ 使\eqref{align:ch07-7.165}成立。此外，对任意凸函数 $E$、满足 $F'=E'\phi'$ 的函数 $F$ 及如\cref{definition:ch07-Pn熵解}的 $\zeta$，有
\begin{align*}
\iint_S\{E(y_n^{\epsilon_j})\partial_t\zeta+F(y_n^{\epsilon_j})\partial_x\zeta\}\,\mathrm{d}x\,\mathrm{d}t
+\int_\mathbb{R}E(u_{0n}^{\epsilon_j})\zeta(x,0)\,\mathrm{d}x
\geqslant\epsilon_j\iint_S F'(y_n^{\epsilon_j})\partial_{xy}y_n^{\epsilon_j}\partial_x\zeta\,\mathrm{d}x\,\mathrm{d}t.
\end{align*}
令 $j\to\infty$ 并如\cref{proposition:ch07-7.4.7}证明中那样论证，得
\begin{align*}
\iint_S[E(y_n)\partial_t\zeta+F(y_n)\partial_x\zeta]\,\mathrm{d}x\,\mathrm{d}t
\geqslant-\int_\mathbb{R}E(u_{0n})\zeta(x,0)\,\mathrm{d}x.
\end{align*}
故 $y_n$ 是 (Pn) 的熵解，由 Kružkov 唯一性定理得 $y_n=u_n$。结论得证。
\end{proof}

\begin{lemma}\label{lemma:ch07-7.6.2}
设 (A0) 与 (A2') 成立，$\phi(0)=0$，$\phi''<0$、$\phi'>0$ 在 $(0,\infty)$。则对所有 $t\in(0,T)$、$\epsilon>0$ 与 $n\in\mathbb{N}$，有
\begin{align}
\pm\frac{\partial}{\partial t}\left[\frac{\alpha\phi(y_n^\epsilon(\cdot,t))+\beta}{t^\alpha}\right]\leqslant0\quad\text{在 }\mathbb{R}\text{ 若 }\pm\alpha>0, \label{align:ch07-7.166a}
\end{align}
\begin{align}
\frac{\partial}{\partial t}[\phi(y_n^\epsilon(\cdot,t))-\beta\log t]\leqslant0\quad\text{在 }\mathbb{R}\text{ 若 }\alpha=0. \label{align:ch07-7.166b}
\end{align}
\end{lemma}

\begin{proof}
为方便置 $A:=\epsilon\partial_{xx}-\partial_x$，于是 $\partial_t y_n^\epsilon=A[\phi(y_n^\epsilon)]$ 在 $S$。令
\begin{align*}
z_n^\epsilon:=t\partial_t y_n^\epsilon-g(y_n^\epsilon),\qquad
g(y_n^\epsilon):=\frac{\alpha\phi(y_n^\epsilon)+\beta}{\phi'(y_n^\epsilon)}\quad(n\in\mathbb{N}).
\end{align*}
由 (A2') 与直接计算得
\begin{align*}
\partial_t z_n^\epsilon=A[\phi'(y_n^\epsilon)z_n^\epsilon]
+\frac{[\alpha+1-g'(y_n^\epsilon)][z_n^\epsilon+g(y_n^\epsilon)]}{t}
\leqslant A[\phi'(y_n^\epsilon)z_n^\epsilon]+\frac{[\alpha+1-g'(y_n^\epsilon)]z_n^\epsilon}{t}
\quad\text{在 }S.
\end{align*}
因 $z_n^\epsilon(\cdot,0)=-g(u_{0n}^\epsilon)\leqslant0$ 在 $\mathbb{R}$，由比较原理得 $z_n^\epsilon\leqslant0$ 在 $S$。从而 $t\partial_t y_n^\epsilon(\cdot,t)\leqslant g(y_n^\epsilon)(\cdot,t)$ 在 $\mathbb{R}$ 对所有 $t\in(0,T)$，由此得\eqref{align:ch07-7.166a}。
\end{proof}

为证\cref{theorem:ch07-7.3.3}(i)，还需要一个引理。

\begin{lemma}\label{lemma:ch07-7.6.3}
设 (A0)--(A1) 成立，$u$ 为所构造解。设 $\{u_{n_k}\}$ 为\cref{theorem:ch07-7.5.1}证明中所用的序列。则对几乎所有 $t\in(0,T)$ 与每个 $x_0\in\operatorname{supp}u_s(\cdot,t)$，存在序列 $\{x_p\}\subseteq\mathbb{R}$ 与子列 $\{u_{n_p}\}\subseteq\{u_{n_k}\}$，使 $x_p\to x_0$ 且 $u_{n_p}(x_p,t)\to\infty$（当 $p\to\infty$）。
\end{lemma}

\begin{proof}
由\cref{remark:ch07-7.5.1}，存在零测集 $N\subseteq(0,T)$ 与 $\{u_{n_k}\}$ 的子列（不重新标号），使
\begin{align}
u_{n_k}(\cdot,t)\stackrel{*}{\rightharpoonup}u_r(\cdot,t)+u_s(\cdot,t)\quad\text{在 }\mathscr{R}_f(I_L) \label{align:ch07-7.167}
\end{align}
对所有 $t\in(0,T)\setminus N$ 与 $L>0$ 成立（见\eqref{align:ch07-7.107}取 $U(y)=y$ 及\eqref{align:ch07-7.109a}--\eqref{align:ch07-7.109b}）。

设 $x_0\in\operatorname{supp}u_s(\cdot,t)$。则不存在邻域 $I_\delta(x_0)$ 使 $\{u_{n_k}(\cdot,t)\}$ 包含于 $L^\infty(I_\delta(x_0))$ 的有界子集；否则存在 $\{u_{n_k}(\cdot,t)\}$ 的子列（不重新标号）与 $f_t\in L^\infty(I_\delta(x_0))$、$f_t\geqslant0$，使 $u_{n_k}(\cdot,t)\stackrel{*}{\rightharpoonup}f_t$ 在 $L^\infty(I_\delta(x_0))$。然而由\eqref{align:ch07-7.167}，这将蕴含 $u_s(\cdot,t)=0$ 在 $I_\delta(x_0)$，矛盾。

由此对所有 $p\in\mathbb{N}$，$\sup_{n_k}\|u_{n_k}(\cdot,t)\|_{L^\infty(I_{1/p}(x_0))}=\infty$。于是存在 $\{u_{n_{k_p}}\}\subseteq\{u_{n_k}\}$ 与 $\{x_p\}$，使 $x_p\in I_{1/p}(x_0)$ 且 $u_{n_p}(x_p,t)\geqslant p$ 对所有 $p\in\mathbb{N}$。结论得证。
\end{proof}

\begin{proof}[\cref{theorem:ch07-7.3.3}的证明]
如前所述，只就 $\phi''<0$、$\phi'>0$ 在 $(0,\infty)$ 的情形证明。

(i) 设 $\phi$ 有界，$\{u_{n_k}\}$ 如\cref{lemma:ch07-7.6.3}证明中。由\cref{lemma:ch07-7.6.1}，对每个 $n_k$，存在 (V$n_{n_k}^\epsilon$) 解的子列 $\{y_{n_k}^{\epsilon_j}\}$，使对几乎所有 $t\in(0,T)$，
\begin{align}
y_{n_k}^{\epsilon_j}(\cdot,t)\to u_{n_k}(\cdot,t)\quad\text{在 }L^1_{\mathrm{loc}}(\mathbb{R})\ \text{当 }j\to\infty. \label{align:ch07-7.168}
\end{align}
另一方面，由\cref{lemma:ch07-7.6.2}的证明，对任意 $t\in(0,T)$，有
\begin{align}
\epsilon_j\partial_{xx}\phi(y_{n_k}^{\epsilon_j}(\cdot,t))-\partial_x\phi(y_{n_k}^{\epsilon_j}(\cdot,t))
=\partial_t y_{n_k}^{\epsilon_j}(\cdot,t)\leqslant\frac{g(y_{n_k}^{\epsilon_j}(\cdot,t))}{t}\quad\text{在 }\mathbb{R}, \label{align:ch07-7.169}
\end{align}
其中 $g(y):=(\alpha\phi(y)+\beta)/\phi'(y)>0$。

固定 $a,b\in\mathbb{R}$，$a<b$，设 $\rho\in C^1_c(a,b)$，$\rho\geqslant0$。置
\begin{align*}
Z(y):=\int_0^y\frac{[\phi'(z)]^2}{\alpha\phi(z)+\beta}\,\mathrm{d}z.
\end{align*}
则 $Z$ 递增，且 $0\leqslant Z(y)\leqslant\phi'(0)$ 对所有 $y\geqslant0$。将\eqref{align:ch07-7.169}乘以 $\rho/g(y_{n_k}^{\epsilon_j}(\cdot,t))$ 并积分，得
\begin{align*}
\int_a^b Z(y_{n_k}^{\epsilon_j})(x,t)[\epsilon_j\rho''(x)+\rho'(x)]\,\mathrm{d}x
\leqslant\frac1t\int_a^b\rho(x)\,\mathrm{d}x
-\epsilon_j\int_a^b\frac{\phi'(y_{n_k}^{\epsilon_j})g'(y_{n_k}^{\epsilon_j})}{[g(y_{n_k}^{\epsilon_j})]^2}[\partial_x y_{n_k}^{\epsilon_j}]^2(x,t)\rho(x)\,\mathrm{d}x
\leqslant\frac1t\int_a^b\rho(x)\,\mathrm{d}x.
\end{align*}
由\eqref{align:ch07-7.168}，令 $j\to\infty$，得
\begin{align}
\int_a^b Z(u_{n_k}(x,t))\rho'(x)\,\mathrm{d}x\leqslant\frac1t\int_a^b\rho(x)\,\mathrm{d}x. \label{align:ch07-7.170}
\end{align}
对几乎所有 $t\in(0,T)$，设 $x_0\in\operatorname{supp}u_s(\cdot,t)$。证明
\begin{align}
\operatorname*{ess\,lim}_{x\to x_0^+}u_r(x,t)=\infty. \label{align:ch07-7.171}
\end{align}
设 $\{x_p\}$ 与 $\{u_{n_p}\}\subseteq\{u_{n_k}\}$ 由\cref{lemma:ch07-7.6.3}给出，固定 $b>x_0$。因 $x_p\to x_0$（当 $p\to\infty$），存在 $\bar p\in\mathbb{N}$ 使 $b>x_p$ 对所有 $p>\bar p$。对任意此种 $p$ 与充分大的 $m\in\mathbb{N}$，置
\begin{align*}
\rho_{p,m}(x):=m(x-x_p)\chi_{[x_p,x_p+1/m]}(x)+\chi_{(x_p+1/m,b-1/m)}(x)+m(b-x)\chi_{[b-1/m,b]}(x).
\end{align*}
在\eqref{align:ch07-7.170}中取 $a=x_p$、$\rho=\rho_{p,m}$，得
\begin{align*}
m\int_{x_p}^{x_p+1/m}Z(u_{n_p}(x,t))\,\mathrm{d}x-m\int_{b-1/m}^b Z(u_{n_p}(x,t))\,\mathrm{d}x
\leqslant\frac1t\int_{x_p}^b\rho_{p,m}(x)\,\mathrm{d}x,
\end{align*}
从而令 $m\to\infty$，对所有 $p>\bar p$ 有
\begin{align*}
Z(u_{n_p}(x_p,t))\leqslant Z(u_{n_p}(b,t))+\frac{b-x_p}{t}.
\end{align*}
由\cref{lemma:ch07-7.6.3}证明，$u_{n_p}(x_p,t)\geqslant p$。令 $p\to\infty$ 并利用\eqref{align:ch07-7.143}，对几乎所有 $b>x_0$，得
\begin{align}
\lim_{y\to\infty}Z(y)\leqslant Z(u_r(b,t))+\frac{b-x_0}{t}, \label{align:ch07-7.172}
\end{align}
从而
\begin{align*}
\lim_{y\to\infty}Z(y)\leqslant\operatorname*{ess\,lim}_{b\to x_0^+}Z(u_r(b,t)).
\end{align*}
因 $Z$ 连续且递增，由上式得\eqref{align:ch07-7.171}。

现可证 $u$ 在 $[0,\tau]$ 于 $x_0\in\operatorname{supp}u_s(\cdot,\tau)$ 处满足相容条件。关于 (C$^-$)，由 $\phi$ 的递增性：事实上若 $u_r>l$ 则 $H_-(u_r-l)=0$；若 $u_r\leqslant l$，则
\begin{align*}
H_-(u_r(x,t)-l)[\phi(u_r(x,t))-\phi(l)]=\phi(l)-\phi(u_r(x,t))\geqslant0,
\end{align*}
故 (C$^-$) 满足。关于 (C$^+$)，由等式\eqref{align:ch07-7.171}，对所有 $l\in[0,\infty)$，对几乎所有 $t\in(0,\tau)$ 有 $\operatorname*{ess\,lim}_{x\to x_0^+}H_-(u_r(\cdot,t)-l)=0$（回忆由单调性此类 $t$ 有 $x_0\in\operatorname{supp}u_s(\cdot,t)$）。于是由控制收敛定理得 (C$^+$)。

(ii) 设 $\phi$ 无界。由\eqref{align:ch07-7.166a}，对所有 $0<t_1\leqslant t\leqslant T$ 与 $x\in\mathbb{R}$，
\begin{align*}
\int_{t_1}^t\phi(y_{n_k}^{\epsilon_j})(x,s)\,\mathrm{d}s
\geqslant
\begin{cases}
\dfrac{\alpha\phi(y_{n_k}^{\epsilon_j})(x,t)+\beta}{\alpha t^\alpha}\dfrac{t^{\alpha+1}-t_1^{\alpha+1}}{\alpha+1}-\dfrac{\beta}{\alpha}(t-t_1)&\text{若 }\alpha\neq0,\\[2mm]
[\phi(y_{n_k}^{\epsilon_j})(x,t)-\beta\log t](t-t_1)+\beta[t\log t-t]-\beta[t_1\log t_1-t_1]&\text{若 }\alpha=0,
\end{cases}
\end{align*}
其中 $\{y_{n_k}^{\epsilon_j}\}$ 为\cref{lemma:ch07-7.6.1}证明中所用子列。令 $\epsilon_j\to0$，由\eqref{align:ch07-7.165}，对几乎所有 $t\in(t_1,T)$ 与 $x\in\mathbb{R}$ 得
\begin{align}
\int_{t_1}^t\phi(u_{n_k}(x,s))\,\mathrm{d}s\geqslant
\begin{cases}
\dfrac{\alpha\phi(u_{n_k})(x,t)+\beta}{\alpha t^\alpha}\dfrac{t^{\alpha+1}-t_1^{\alpha+1}}{\alpha+1}-\dfrac{\beta}{\alpha}(t-t_1)&\text{若 }\alpha\neq0,\\
[\phi(u_{n_k})(x,t)-\beta](t-t_1)+\beta t_1\log\dfrac{t}{t_1}&\text{若 }\alpha=0.
\end{cases} \label{align:ch07-7.173}
\end{align}
另一方面，如\cref{proposition:ch07-7.2.1}证明中那样，由弱形式\eqref{align:ch07-7.78}易见，对每个 $t\in(0,T]$，映射 $x\mapsto Y_n(x):=\int_0^t\phi(u_{n_k}(x,s))\,\mathrm{d}s$ 属于 $W^{1,1}(\mathbb{R})$，其弱导数 $Y_n'(\cdot)=u_{0n}-u_{n_k}(\cdot,t)$ 在 $\mathbb{R}$ 中几乎处处。故
\begin{align*}
0\leqslant Y_n\leqslant\|u_{0n}\|_{L^1(\mathbb{R})}\leqslant\|u_0\|_{\mathscr{R}_f(\mathbb{R})}.
\end{align*}
结合此不等式与\eqref{align:ch07-7.173}并令 $t_1\to0^+$，得
\begin{align}
\phi(u_{n_k})(x,t)\leqslant\frac{(\alpha+1)\|u_0\|_{\mathscr{R}_f(\mathbb{R})}}{t}+\beta
\quad\text{对几乎所有 }t\in(0,T)\text{ 与 }x\in\mathbb{R} \label{align:ch07-7.174}
\end{align}
（回忆对无界 $\phi$，因假设 $\phi'>0$ 且 $\alpha\phi+\beta>0$ 在 $[0,\infty)$，故 $\alpha\geqslant0$）。由\eqref{align:ch07-7.174}，序列 $\{u_{n_k}\}$ 对每个 $\tau\in(0,T)$ 包含于 $L^\infty(\mathbb{R}\times(\tau,T))$ 的有界子集，于是由\cref{theorem:ch07-7.3.1}，令 $n\to\infty$ 即得结论。
\end{proof}

\begin{proof}[\cref{proposition:ch07-7.3.5}的证明]
由\cref{theorem:ch07-7.3.3}的证明，对几乎所有 $t\in(0,T)$，不等式\eqref{align:ch07-7.172}对所有 $x_j\in\operatorname{supp}u_s(\cdot,t)$ 成立。设 $x_1\in\operatorname{supp}u_s(\cdot,t)$，置 $I_1:=(x_1-\epsilon,x_1+\epsilon)$（$\epsilon>0$）。因 $Z$ 可逆，由\eqref{align:ch07-7.172}得
\begin{align*}
\int_{I_1}u_r(x,t)\,\mathrm{d}x\geqslant\int_{x_1}^{x_1+\epsilon}Z^{-1}\!\left(\lim_{y\to\infty}Z(y)-\frac{x-x_1}{t}\right)\mathrm{d}x
=\int_0^\epsilon Z^{-1}\!\left(\lim_{y\to\infty}Z(y)-\frac{y}{t}\right)\mathrm{d}y=:B_\epsilon,
\end{align*}
其中
\begin{align}
\lim_{\epsilon\to0^+}\frac{B_\epsilon}{\epsilon}=Z^{-1}\!\left(\lim_{y\to\infty}Z(y)\right)=\infty. \label{align:ch07-7.175}
\end{align}
若 $\operatorname{supp}u_s(\cdot,t)\not\subseteq I_1$，取 $x_2\in\operatorname{supp}u_s(\cdot,t)\setminus I_1$，置 $I_2:=(x_2-\epsilon,x_2+\epsilon)$。因 $(x_1,x_1+\epsilon)\cap(x_2,x_2+\epsilon)=\emptyset$，
\begin{align*}
\int_{I_1\cup I_2}u_r(x,t)\,\mathrm{d}x\geqslant\int_{x_1}^{x_1+\epsilon}u_r\,\mathrm{d}x+\int_{x_2}^{x_2+\epsilon}u_r\,\mathrm{d}x\geqslant2B_\epsilon.
\end{align*}
递归地，只要 $\operatorname{supp}u_s(\cdot,t)\not\subseteq\bigcup_{i=1}^{n-1}I_i$（$I_i:=(x_i-\epsilon,x_i+\epsilon)$），就存在 $x_n\in\operatorname{supp}u_s(\cdot,t)\setminus\bigcup_{i=1}^{n-1}I_i$，使
\begin{align*}
nB_\epsilon\leqslant\int_{\bigcup_{i=1}^n I_i}u_r(x,t)\,\mathrm{d}x\leqslant\|u_0\|_{\mathscr{R}_f(\mathbb{R})}.
\end{align*}
故过程在某个 $n=n_\epsilon$ 处终止，$n_\epsilon B_\epsilon\leqslant\|u_0\|_{\mathscr{R}_f(\mathbb{R})}$。于是
\begin{align*}
\operatorname{supp}u_s(\cdot,t)\subseteq\bigcup_{i=1}^{n_\epsilon}I_i,
\end{align*}
从而
\begin{align*}
\lambda(\operatorname{supp}u_s(\cdot,t))\leqslant\lambda\!\left(\bigcup_{i=1}^{n_\epsilon}I_i\right)\leqslant2n_\epsilon\epsilon
\leqslant\frac{2\|u_0\|_{\mathscr{R}_f(\mathbb{R})}\epsilon}{B_\epsilon}.
\end{align*}
令 $\epsilon\to0^+$ 并利用\eqref{align:ch07-7.175}，即得结论。
\end{proof}

\end{document}
